\documentclass[11pt]{article}

\usepackage[margin=1in]{geometry}
\usepackage{amsmath,amssymb,amsfonts,amsthm}
\usepackage{array}
\usepackage{bm}
\usepackage{graphicx}
\usepackage{multirow}
\usepackage{multicol}
\usepackage{nicefrac}
\usepackage{paralist}
\usepackage{enumitem}
\usepackage{verbatim}
\usepackage{thm-restate}
\usepackage[normalem]{ulem}
\usepackage[compact]{titlesec}
\usepackage{tikz}
\usetikzlibrary{arrows}
\usetikzlibrary{arrows.meta}
\usetikzlibrary{backgrounds}
\usetikzlibrary{calc}
\usetikzlibrary{decorations.markings}
\usetikzlibrary{fit}
\usetikzlibrary{patterns}
\usetikzlibrary{positioning}
\usetikzlibrary{shapes}

\usepackage{fancybox}
\newenvironment{fminipage}%
  {\begin{Sbox}\begin{minipage}}%
  {\end{minipage}\end{Sbox}\fbox{\TheSbox}}

\newenvironment{algbox}[0]{%
\vskip 0.2in
\noindent
\begin{fminipage}{6.3in}
}{%
\end{fminipage}
\vskip 0.2in
}

\usepackage{tcolorbox}
\tcbuselibrary{skins,breakable}
\tcbset{enhanced jigsaw}
\usepackage{mdframed}
\usepackage{xcolor}
\usepackage{cite}
\usepackage{url}
\usepackage{xspace}
\usepackage[mathscr]{euscript}
\usepackage{nameref}
\usepackage[
  linktocpage=true,
  pagebackref=true,
  colorlinks=true,
  linkcolor=blue,
  citecolor=blue,
  urlcolor=blue,
  bookmarks,
  bookmarksopen,
  bookmarksnumbered
]{hyperref}
\usepackage[noabbrev,nameinlink]{cleveref}

\newtheorem{theorem}{Theorem}[section]
\newtheorem{example}{Example}[section]
\newtheorem{lemma}{Lemma}[section]
\newtheorem{proposition}[lemma]{Proposition}
\newtheorem{corollary}[lemma]{Corollary}
\newtheorem{claim}[lemma]{Claim}
\newtheorem{fact}[lemma]{Fact}
\newtheorem{conjecture}[lemma]{Conjecture}
\newtheorem{definition}[lemma]{Definition}
\newtheorem{problem}{Problem}
\newtheorem{remark}[lemma]{Remark}
\newtheorem{observation}[lemma]{Observation}

\newtheorem*{theorem*}{Theorem}
\newtheorem*{claim*}{Claim}
\newtheorem*{proposition*}{Proposition}
\newtheorem*{lemma*}{Lemma}
\newtheorem*{problem*}{Problem}

\allowdisplaybreaks
\setlength{\parskip}{3pt}

\title{Convergence of Normalized Residuals for Restarted Conjugate Gradients of Length Two}
\author{}
\date{July 29, 2026}

\begin{document}

\pagenumbering{roman}

\maketitle

\begin{abstract}
Let \(A\in\mathbb R^{n\times n}\) be symmetric positive definite, let
\(f(x)=\frac12x^TAx-b^Tx\), and set \(r_k=b-Ax_k\).  In exact arithmetic,
define each restart of length two by minimizing \(f\) over
\(x_k+\operatorname{span}\{r_k,Ar_k\}\).  If the initial residual has
nonzero spectral projection on at least three distinct eigenspaces of
\(A\), then every residual remains nonzero with spectral grade at least
three, and the even and odd subsequences of Euclidean-normalized residuals
converge to unit vectors \(y_{\mathrm e}\) and \(y_{\mathrm o}\),
respectively.  Both limits have spectral grade three or four and form a
two-cycle under the normalized monic-quadratic residual map.  Moreover, if
\(E_k=\lVert r_k\rVert_{A^{-1}}^2\), then
\[
  0<\frac{E_{k+1}}{E_k}\leq\frac{E_{k+2}}{E_{k+1}}
  \leq
  \left(\frac{\lambda_{\max}-\lambda_{\min}}
  {\lambda_{\max}+\lambda_{\min}}\right)^2<1,
\]
where \(\lambda_{\min}\) and \(\lambda_{\max}\) are the extreme eigenvalues
in the initial cyclic subspace.  No simple-spectrum or genericity
assumption is imposed.  Each restart uses two matrix--vector products and
\(O(n)\) additional arithmetic operations.
\end{abstract}

\noindent\textbf{Note.} This paper was fully AI generated by a harness wrapped around GPT 5.6 Sol Ultra from a one-shot prompt by Richard Peng.

\clearpage
\tableofcontents
\clearpage

\pagenumbering{arabic}

\section{Introduction}
\label{sec:introduction}

Conjugate gradients (CG) solves a symmetric positive definite linear system
\(Ax=b\) by minimizing
\[
 f(x)=\frac12x^TAx-b^Tx
\]
over expanding Krylov spaces.  In exact arithmetic, the unrestarted method
terminates after at most the number of distinct eigenvalues active in the
initial residual \cite{HestenesStiefel1952}.  A fixed restart length changes
the question: after \(s\) inner CG steps, the accumulated conjugacy
relations are discarded and the same construction begins again from the
new residual.  Thus, if \(r_k=b-Ax_k\), one outer step minimizes \(f\) over
\[
 x_k+\mathcal K_s(A,r_k),\qquad
 \mathcal K_s(A,r_k)
 =\operatorname{span}\{r_k,Ar_k,\ldots,A^{s-1}r_k\}.
\]
Forsythe called this the optimum \(s\)-gradient method; in modern terms it
is CG restarted every \(s\) steps
\cite{Forsythe1968,FaberLiesenTichy2023}.

Even as the iterates approach the solution and the residuals tend to zero,
their normalized directions may oscillate.
For restart length \(s=1\), the iteration is exact-line-search steepest
descent.  Forsythe and Motzkin formulated the corresponding direction
problem, and Akaike proved that the even and odd normalized residuals
converge separately to two alternating directions supported on the extreme
active eigenspaces
\cite{ForsytheMotzkin1951,Akaike1959}.  Forsythe's general conjecture asks
for the same parity convergence at every fixed restart length: each of
\(\{r_{2k}/\lVert r_{2k}\rVert\}\) and
\(\{r_{2k+1}/\lVert r_{2k+1}\rVert\}\) should have a single limit
\cite{Forsythe1968,FaberLiesenTichy2023}.

We prove this statement for restart length two.  For \(v\neq0\), let
\(d(A,v)\) denote its spectral grade, namely the number of distinct
eigenvalues of \(A\) on whose eigenspaces \(v\) has nonzero projection.
For a subspace \(V\), write \(u\perp_A V\) when
\(\langle Au,v\rangle=0\) for every \(v\in V\).

\begin{theorem}[Forsythe conjecture for restart length two]
\label{thm:forsythe-s2}
In exact arithmetic, let \(A\in\mathbb R^{n\times n}\) be symmetric
positive definite, let \(b,x_0\in\mathbb R^n\), put
\(x_\star=A^{-1}b\) and \(r_0=b-Ax_0\), and assume
\(d(A,r_0)\geq3\).  For \(k\geq0\), define
\[
 \mathcal K_2(A,r_k)=\operatorname{span}\{r_k,Ar_k\},\qquad
 r_k=b-Ax_k,
\]
and let \(x_{k+1}\) be the unique vector satisfying
\[
 x_{k+1}\in x_k+\mathcal K_2(A,r_k),\qquad
 x_\star-x_{k+1}\perp_A\mathcal K_2(A,r_k).
\]
Equivalently, these are the iterates generated by
\hyperref[alg:restarted-cg2]{\textsc{Restarted-CG(2)}}.  Then
\(d(A,r_k)\geq3\) for every \(k\).  If
\(E_k=\lVert r_k\rVert_{A^{-1}}^2\), and if
\(\lambda_{\min}\) and \(\lambda_{\max}\) are the smallest and largest
eigenvalues on which \(r_0\) has nonzero spectral projection, then
\[
 0<\frac{E_{k+1}}{E_k}\leq\frac{E_{k+2}}{E_{k+1}}
 \leq
 \left(\frac{\lambda_{\max}-\lambda_{\min}}
 {\lambda_{\max}+\lambda_{\min}}\right)^2<1.
\]
Furthermore, there are unit vectors \(y_{\mathrm e},y_{\mathrm o}\) such
that
\[
 \frac{r_{2k}}{\lVert r_{2k}\rVert}\longrightarrow y_{\mathrm e},
 \qquad
 \frac{r_{2k+1}}{\lVert r_{2k+1}\rVert}\longrightarrow y_{\mathrm o}.
\]
Both limit vectors have grade three or four.  For the monic-quadratic map
\(T\) defined in \Cref{eq:monic-map}, they satisfy
\(y_{\mathrm o}=T(y_{\mathrm e})\) and
\(y_{\mathrm e}=T(y_{\mathrm o})\).
The deterministic implementation produces the iterates, residuals, and
normalized residuals.  If one multiplication by \(A\) costs \(T_A\), then
each restart costs \(2T_A+O(n)\) arithmetic operations, and hence
\(O(n^2)\) operations for an explicit dense matrix.
\end{theorem}

The grade hypothesis is the natural nontermination threshold.  A residual
of grade at most two is annihilated by a CG residual polynomial of degree
at most two, after which its direction is undefined.  The theorem shows
that grade at least three instead persists at every restart.  It also
aggregates repeated eigenvalues through their spectral projections, so no
simple-spectrum assumption is needed, and it imposes no genericity
condition on the starting residual.  The two limiting directions are
companion points under \(T\), and their grades are rigidly restricted to
three or four.

\paragraph{Results for the same iteration.}
The classical \(s=1\) result and the case
\(d(A,r_0)=s+1\), in which Forsythe showed that the normalized process is
already an exact two-cycle, are restricted versions of the same problem
\cite{Akaike1959,Forsythe1968}.  For \(s=2\), Zhuk and Bondarenko
subsequently presented a confirmation of Forsythe's conjecture and derived
convergence of the relevant polynomial coefficients together with rate and
asymptotic-spectrum results \cite{ZhukBondarenko1984}.  Faber, Liesen, and
Tich\'y later observed that the passage from coefficient convergence to a
single residual direction relies on an implication attributed to
Zabolotskaya; the accessible English translation describes the coefficient
limit as indicating, rather than establishing, a unique direction
\cite{Zabolotskaya1979,FaberLiesenTichy2023}.  The present proof supplies
the direction-convergence argument directly, without treating coefficient
convergence alone as sufficient.

Pronzato, Wynn, and Zhigljavsky studied exactly the same optimum
\(s\)-gradient iteration as a dynamical system on spectral probability
measures \cite{PronzatoWynnZhigljavsky2009}.  Their \(s=2\) analysis
establishes a singleton in the three-point regime, while fixed limiting
invariants can leave a one-parameter family of four-point candidates.
Faber, Liesen, and Tich\'y obtained a modern projection formulation,
two-step asymptotic regularity, connected parity limit sets, and the bound
\(s<d\leq2s\) for every limiting direction
\cite{FaberLiesenTichy2023}.  At \(s=2\), these results isolate grades three
and four but do not turn a possible four-point continuum into a single
signed direction.  \Cref{thm:forsythe-s2} covers both grades for every
finite-dimensional instance satisfying the stated nontermination
hypothesis.

\paragraph{Nearby but different iterations.}
The nonsymmetric Arnoldi cross iteration of
\cite{FaberLiesenTichy2023} alternates projection steps for \(A\) and
\(A^T\); only its symmetric specialization has the direction dynamics
studied here.  Restarted GMRES instead minimizes the Euclidean residual
norm, and its cycle-convergence results concern a different process
\cite{SaadSchultz1986,VecharynskiLangou2010}.  Likewise, the
communication-avoiding ``\(s\)-step CG'' method retains conjugacy across
blocks and is not CG restarted after every \(s\) inner steps
\cite{ChronopoulosGear1989}.  These methods provide neighboring Krylov
context, but they are not same-problem partial results.

\paragraph{Proof perspective and organization.}
The proof represents one restart by a monic quadratic that is orthogonal
for the current spectral measure.  Two monotone scalar quantities then
force a limiting two-restart quartic and a connected omega-limit set
supported on three or four spectral coordinates.  Three coordinates reduce
to a finite algebraic classification.  Four coordinates admit a continuum
of exact two-cycles, so polynomial convergence alone does not rule out
motion along that continuum.  A corrected barycentric coordinate retains
the effect of vanishing outside coordinates, while a Lyapunov-tail estimate
and a logarithmic cocycle exclude the remaining resonant drift.

\Cref{sec:preliminaries} fixes the linear-algebra, spectral, and
orthogonal-polynomial notation.  \Cref{sec:overview} gives the proof
architecture and the main identities.  \Cref{sec:main-proof} proves the
theorem, beginning with the exact restarted-CG implementation and ending
with the collapse of the three- and four-coordinate omega-limit sets.

\section{Preliminaries}
\label{sec:preliminaries}

\paragraph{Linear-algebra conventions.}
All arithmetic is exact.  Inner products and norms without a subscript are
Euclidean.  For a symmetric positive definite matrix \(A\), write
\[
 \langle u,v\rangle_A=\langle Au,v\rangle,\qquad
 \langle u,v\rangle_{A^{-1}}=\langle A^{-1}u,v\rangle
\]
for the energy and inverse-energy inner products, with corresponding norms
\(\lVert\cdot\rVert_A\) and \(\lVert\cdot\rVert_{A^{-1}}\).  The matrix
norm \(\lVert A\rVert\) is the induced Euclidean operator norm.  For a
subspace \(V\), the notation \(u\perp_A V\) means
\(\langle u,v\rangle_A=0\) for every \(v\in V\).

For fixed \(A\) and \(b\), put
\[
 x_\star=A^{-1}b,\qquad
 f(x)=\frac12x^TAx-b^Tx,\qquad
 r(x)=b-Ax=A(x_\star-x),
\]
and set \(\mathcal K_2(A,r)=\operatorname{span}\{r,Ar\}\).  The standard
quadratic identities
\[
 \nabla f(x)=-r(x),\qquad
 f(x)-f(x_\star)
 =\frac12\lVert x-x_\star\rVert_A^2
 =\frac12\lVert r(x)\rVert_{A^{-1}}^2
\]
show that, for any subspace \(V\) and any \(x^+\in x+V\),
\[
 x^+=\operatorname*{argmin}_{z\in x+V}f(z)
 \quad\Longleftrightarrow\quad
 r(x^+)\perp V
 \quad\Longleftrightarrow\quad
 x_\star-x^+\perp_A V.
\]
For a nonzero residual \(r_k=r(x_k)\), we write
\(y_k=r_k/\lVert r_k\rVert\).

\paragraph{Spectral data.}
\begin{definition}[Spectral grade and support]
\label{def:grade-support}
Let \(\Pi_\lambda\) be the orthogonal spectral projector of a symmetric
matrix \(A\) associated with \(\lambda\).  For \(v\in\mathbb R^n\), define
\[
 \Lambda(A,v)=\{\lambda:\Pi_\lambda v\neq0\},\qquad
 d(A,v)=\lvert\Lambda(A,v)\rvert.
\]
For a coordinate vector \(z\), write
\(\operatorname{supp}(z)=\{i:z_i\neq0\}\).
\end{definition}

For \(v\neq0\), the cyclic subspace generated by \(v\) is
\[
 \operatorname{span}\{p(A)v:p\in\mathbb R[t]\}
 =\operatorname{span}\{\Pi_\lambda v:\lambda\in\Lambda(A,v)\}.
\]
The spectral theorem gives
\[
 p(A)v=\sum_{\lambda\in\Lambda(A,v)}
             p(\lambda)\Pi_\lambda v.
\]
Thus \(p(A)v=0\) exactly when \(p\) vanishes on
\(\Lambda(A,v)\).  When \(A\) and \(v\) are clear, the extreme elements
of this active spectrum are denoted by
\(\lambda_{\min}\) and \(\lambda_{\max}\).

For a unit vector \(y\), its spectral measure and moments are
\[
 \nu_y=\sum_{\lambda\in\Lambda(A,y)}
          \lVert\Pi_\lambda y\rVert^2\delta_\lambda,\qquad
 \mu_j(y)=\int t^j\,d\nu_y(t)
         =\langle A^jy,y\rangle\quad(j\geq0).
\]
Here \(\delta_\lambda\) is the unit point mass at \(\lambda\).

\begin{definition}[Normalized monic-quadratic map]
\label{def:monic-map}
For a unit vector \(y\) with \(d(A,y)\geq3\), let
\[
 \pi_y(t)=t^2+c_1(y)t+c_0(y)
\]
be the monic quadratic satisfying
\[
 \langle\pi_y(A)y,y\rangle=0,\qquad
 \langle\pi_y(A)y,Ay\rangle=0.
\]
Its coefficients are characterized by
\[
 \begin{pmatrix}
   \mu_1(y)&\mu_0(y)\\
   \mu_2(y)&\mu_1(y)
 \end{pmatrix}
 \begin{pmatrix}c_1(y)\\c_0(y)\end{pmatrix}
 =
 -\begin{pmatrix}\mu_2(y)\\\mu_3(y)\end{pmatrix}.
\]
Strict Cauchy--Schwarz gives
\[
 c_0(y)=
 \frac{\mu_1(y)\mu_3(y)-\mu_2(y)^2}
      {\mu_0(y)\mu_2(y)-\mu_1(y)^2}>0.
\]
Define
\[
 \theta(y)=\lVert\pi_y(A)y\rVert,\qquad
 T(y)=\frac{\pi_y(A)y}{\theta(y)},\qquad
 P_y(t)=\frac{\pi_y(t)}{c_0(y)}.
\tag{1}\label{eq:monic-map}
\]
\end{definition}

A quadratic cannot vanish on three distinct active eigenvalues, so
\(\theta(y)>0\).  Hence the objects in
\Cref{def:monic-map} are well defined, \(P_y(0)=1\), and
\(c_1,c_0,\theta,T\), and \(y\mapsto P_y\) are continuous on the
grade-at-least-three domain.
For every monic quadratic \(q\), orthogonality and the fact that
\(q-\pi_y\) is linear give
\[
 \lVert q(A)y\rVert^2
 =\lVert\pi_y(A)y\rVert^2+
   \lVert(q-\pi_y)(A)y\rVert^2.
\tag{2}\label{eq:monic-minimization}
\]
Thus \(\theta(y)\) is the minimum of
\(\lVert q(A)y\rVert\) over monic quadratics \(q\).

\paragraph{Limit sets.}
For a precompact sequence \(\{z_k\}\), let
\(\omega(z_k)\) denote its set of subsequential limits.  This set is
nonempty and compact.  If
\(\lVert z_{k+1}-z_k\rVert\to0\), then \(\omega(z_k)\) is connected;
if \(\omega(z_k)\) is a singleton, then \(z_k\) converges.

\section{Overview}
\label{sec:overview}

The proof converts restarted conjugate gradients into a dynamical system on
the unit sphere and then shows that the even omega-limit set of this system
is a singleton.  The common part of the argument produces a rigid limiting
quartic whose equality set contains only three or four active eigenvalues.
The grade-three case then collapses by finite-dimensional algebra.  The
main technical challenge is the four-eigenvalue case: the exact dynamics on
a four-node spectral face has a continuum of two-cycles, so convergence of
the quartic does not by itself prevent drift along that continuum.  Moreover,
spectral coordinates outside the face may vanish without geometric decay.
We control this drift by a corrected barycentric coordinate, charge the
exceptional transverse mass to a Lyapunov tail, and exclude the remaining
multiplier-\(-1\) resonance with a logarithmic cocycle.

\paragraph{Polynomial and spectral reduction.}
In exact arithmetic, the algorithm
\hyperref[alg:restarted-cg2]{\textsc{Restarted-CG(2)}} takes
\(A,b,x_0\), with \(A\) symmetric positive definite,
\(r_0=b-Ax_0\), and \(d(A,r_0)\geq3\), and returns the iterates \(x_k\),
residuals \(r_k=b-Ax_k\), and directions
\(y_k=r_k/\lVert r_k\rVert\).
At each restart it minimizes the quadratic objective over
\(x_k+\operatorname{span}\{r_k,Ar_k\}\), or equivalently makes
\(r_{k+1}\) orthogonal to that Krylov space.  The two optimality equations
give the exact residual-polynomial form
\[
 r_{k+1}=P_{y_k}(A)r_k,\qquad
 P_y(t)=\frac{\pi_y(t)}{c_0(y)},\qquad
 y_{k+1}=T(y_k)=\frac{\pi_{y_k}(A)y_k}{\theta(y_k)}.
\]
Here \(\pi_y\), \(c_0(y)=\pi_y(0)>0\), and \(\theta(y)\) are as in
\Cref{def:monic-map}.  The actual residual polynomial is normalized to have
constant term one, while its monic multiple \(\pi_y\) is useful because it
is the minimum-norm monic quadratic in \Cref{eq:monic-minimization}; its
norm will be the first Lyapunov quantity.  A fixed spectral isometry then
replaces each active eigenspace by one coordinate, reducing the problem to
\(A=\operatorname{diag}(\lambda_1,\ldots,\lambda_m)\) with distinct
eigenvalues.  This loses no information and is why repeated eigenvalues
require no genericity assumption.  Computationally, one restart forms
\(Ar_k,A^2r_k\), solves a \(2\times2\) system, and performs vector updates,
for a cost of \(2T_A+O(n)\) when one multiplication by \(A\) costs
\(T_A\), or \(O(n^2)\) for a dense matrix; see
\Cref{lem:iteration-equivalence,lem:spectral-reduction}.

\paragraph{Two scalar limits.}
Put \(\theta_k=\theta(y_k)\).  The monic-projection identity gives
\[
 \theta_k=\theta_{k+1}\langle y_k,y_{k+2}\rangle,\qquad
 \lVert y_{k+2}-y_k\rVert^2
 =2\left(1-\frac{\theta_k}{\theta_{k+1}}\right).
\]
Thus \(\theta_k\) is nondecreasing and converges to a number \(\tau>0\),
and the squared two-step displacements are summable.  Independently, if
\(E_k=\lVert r_k\rVert_{A^{-1}}^2\) and
\(\rho_k=E_{k+1}/E_k\), the cross identity
\(\langle r_k,r_{k+2}\rangle_{A^{-1}}=E_{k+1}\) yields
\[
 0<\rho_k\leq\rho_{k+1}\leq
 \left(\frac{\lambda_{\max}-\lambda_{\min}}
 {\lambda_{\max}+\lambda_{\min}}\right)^2<1.
\]
Consequently \(\rho_k\to\rho_\infty\in(0,1)\).  The same argument proves
that no residual vanishes and that grade at least three persists, so the
normalized orbit is defined for all time.  The second scalar limit is also
needed below to fix the constant coefficient of the limiting quartic when
only three interpolation nodes are present; see
\Cref{lem:monic-monotonicity,lem:energy-ratio}.

\paragraph{Omega-limit rigidity and the grade-three case.}
Let \(\Omega=\omega(y_{2k})\).  Compactness of the unit sphere and vanishing
two-step displacement make \(\Omega\) nonempty, compact, and connected.
Continuity of the monic map shows that every \(y\in\Omega\) satisfies
\[
 T^2y=y,\qquad
 q_y(t):=\pi_{Ty}(t)\pi_y(t),\qquad
 q_y(A)y=\beta y,\qquad \beta:=\tau^2.
\]
Thus every active eigenvalue of \(y\) is a root of the monic quartic
\(q_y-\beta\), while \(\tau>0\) rules out grade below three.  Every
omega-limit point therefore has grade three or four.

The energy-ratio limit removes the remaining ambiguity among these
quartics.  In fact,
\[
 q_{2k}:=\pi_{y_{2k+1}}\pi_{y_{2k}}\longrightarrow q_\star,\qquad
 q_\star=\pi_\star^{\mathrm o}\pi_\star^{\mathrm e},\qquad
 q_\star(0)=\frac{\beta}{\rho_\infty},
\]
and the even and odd quadratic factors converge separately.  The reason is
that \(q_y(\lambda)=\beta\) on the support of \(y\), and the displayed
constant coefficient supplies the fourth condition on a three-point
support.  There are only finitely many possible supports, so connectedness
forces one quartic throughout \(\Omega\).  With
\[
 I_\star=\{i:q_\star(\lambda_i)=\beta\},
 \qquad 3\leq\lvert I_\star\rvert\leq4,
 \qquad
 \sum_{k=0}^{\infty}\sum_{j\notin I_\star}y_{2k,j}^2<\infty,
\]
the even-parity mass away from the quartic equality set is summable.  If
every point of \(\Omega\) has grade three, a support and a quadratic divisor of
\(q_\star\) determine the squared coordinates through the two orthogonality
equations and normalization.  Only finitely many supports, divisors, and
coordinate-sign choices remain.  Hence the connected set \(\Omega\) is
finite and therefore a singleton.  This proves the grade-three regime,
including grade-three boundary accumulation when
\(\lvert I_\star\rvert=4\); see
\Cref{lem:omega-structure,lem:quartic-convergence,%
lem:off-support-summability,lem:three-point-collapse}.

\paragraph{Corrected four-node transport.}
It remains to consider an omega-limit point positive on four equality
nodes \(\xi_1<\xi_2<\xi_3<\xi_4\).  Put
\[
 L(t)=\prod_{i=1}^4(t-\xi_i).
\]
Then \(q_\star=L+\beta\).  On the exact four-node face, Lagrange
interpolation supplies a barycentric parameter \(\zeta\); when the cubic
coefficient defect vanishes, two normalized steps return the corresponding
state.  The resulting family of exact two-cycles is the source of the
possible drift.  Coordinates outside \(I_\star\) cannot simply be deleted,
because they may remain nonzero at finite times and enter the orthogonality
equations.

For the full orbit, write \(p_n=\pi_{y_n}\),
\(H_n=\theta_n^2\), \(q_n=p_np_{n+1}\), let \(v_{n,i}\) be the squared
mass at \(\xi_i\), and let \(w_{n,j}=y_{n,j}^2\) outside \(I_\star\).
With
\[
 \ell_i(x)=\frac{L(x)}{L'(\xi_i)(x-\xi_i)},
\]
the corrected coordinate \(\zeta_n\) is defined by the four identities
\[
 v_{n,i}p_n(\xi_i)
 +\sum_{j\notin I_\star}
   w_{n,j}p_n(\lambda_j)\ell_i(\lambda_j)
 =H_n\frac{\xi_i-\zeta_n}{L'(\xi_i)}
 \qquad(1\leq i\leq4).
\]
The Lagrange correction is chosen precisely so that all outside
coordinates are retained and the resulting transport law is exact.  Set
\[
 \delta_n=[t^3]q_n+\sum_{i=1}^4\xi_i,\qquad
 \epsilon_n=\sum_{j\notin I_\star}w_{n,j}.
\]
Here \([t^3]q_n\) denotes the cubic coefficient of \(q_n\).
The outside-coordinate terms collect into an explicit polynomial \(F_n\)
whose coefficients are \(O(\epsilon_n)\), and the main transport identity
becomes
\[
 \zeta_{n+1}-\zeta_n
 =\frac{\delta_nL(\zeta_n)-F_n(\zeta_n)/H_n}{H_{n+1}}.
\]
An accompanying exact defect-transport identity relates
\(\delta_{n+1}\) to \(\delta_n\).  These two identities are the mechanism
that turns decay of the transverse mass into convergence along the
four-node family; see
\Cref{lem:four-node-coordinate,lem:forced-barycentric}.

\paragraph{Tail control, resonance exclusion, and collapse.}
For a surviving outside coordinate \(\lambda_j\), the limiting two-step
multiplier is \(q_\star(\lambda_j)/\beta\).  A modulus greater than one is
incompatible with the already proved decay, \(+\!1\) would put \(j\) in
\(I_\star\), and a strict modulus below one gives geometric decay.  The
only noncontracting possibility is therefore the signed multiplier \(-1\).
For \(\Delta_n=\lVert y_{n+2}-y_n\rVert^2\), the exact Lyapunov tail
\[
 \tau-\theta_n
 =\frac12\sum_{\ell=n}^{\infty}\theta_{\ell+1}\Delta_\ell
\]
charges the mass of every such resonant coordinate.  Together with the
strictly contracting coordinates and the forced transport identities, it
gives constants \(C>0\) and \(0<\gamma<1\) such that
\[
 \sum_{\ell=n}^{\infty}(\Delta_\ell+\epsilon_\ell)
 \leq C(\tau-\theta_n+\gamma^n),\qquad
 \lvert\delta_n\rvert\leq C(\tau-\theta_n+\gamma^n).
\]
Endpoint estimates are essential for the defect bound when \(\zeta_n\)
approaches \(\xi_2\) or \(\xi_3\) and an active squared coordinate tends
to zero.

For a surviving resonant node \(\lambda=\lambda_j\), the proof introduces
the exact logarithmic cocycle
\[
 \Psi_{n,\lambda}
 =\log w_{n,j}+\sum_{i=1}^4\ell_i(\lambda)\log v_{n,i}.
\]
Writing \(e_n=\beta/(\theta_n\theta_{n+1})-1\) and
\(b_{n,i}=q_n(\xi_i)/(\theta_n\theta_{n+1})-1\), cubic interpolation
cancels every term linear in the active errors and yields
\[
 \frac{\Psi_{n+2,\lambda}-\Psi_{n,\lambda}}2
 =2e_n+O\left(e_n^2+\sum_{i=1}^4b_{n,i}^2\right).
\]
The Lyapunov and endpoint bounds make the negative parts of these
increments summable.  Thus the cocycle is bounded below on each parity,
contradicting \(w_{n,j}\to0\) along a subsequence approaching a
four-positive omega-limit point.  No multiplier-\(-1\) coordinate survives.

All remaining outside mass now decays geometrically, and the exact defect
transport reduces to the parabolic estimate
\[
 \lvert\delta_{n+1}-\delta_n\rvert
 \leq C\delta_n^2+C\gamma^n.
\]
The sign-or-summability dichotomy says that \(\delta_n\) is either
summable or eventually one-signed.  In the first case the barycentric
transport gives finite total variation of \(\zeta_n\); in the second,
\(L\geq0\) on \([\xi_2,\xi_3]\) makes the adverse increments summable.
Hence \(\zeta_n\) converges in either case.  The corrected barycentric
identities then recover convergence of the four active squared
coordinates, and the positive limiting two-step multipliers stabilize
their signs.  Thus \(y_{2k}\) converges.  Continuity of \(T\) gives
convergence of \(y_{2k+1}\) to its companion point in the \(T\)-two-cycle,
both limits have grade three or four, and the spectral isometry transfers
the conclusion back to the original residual directions.  This completes
\Cref{thm:forsythe-s2}; see
\Cref{lem:outside-tail,lem:defect-tail,lem:endpoint-errors,%
lem:no-surviving-resonance,lem:strict-gap-collapse,%
lem:four-point-collapse}.

\section{Proof of the Main Theorem}
\label{sec:main-proof}

\subsection{Polynomial form of one restart}

The normalized monic-quadratic map from \Cref{def:monic-map} is the
spectral form of one restarted CG step.  The following algorithm gives its
exact-arithmetic implementation.

\begin{algbox}
\textbf{\uline{\textsc{Restarted-CG(2)}}}: Perform conjugate-gradient
restarts of length two.
\phantomsection\label{alg:restarted-cg2}

\uline{Input}: A symmetric positive definite matrix
\(A\in\mathbb R^{n\times n}\), vectors \(b,x_0\in\mathbb R^n\), and
\(r_0=b-Ax_0\) with \(d(A,r_0)\geq3\).

\uline{Output}: Sequences \(\{x_k\}_{k\geq0}\),
\(\{r_k\}_{k\geq0}\), and
\(\{y_k\}_{k\geq0}\), where \(r_k=b-Ax_k\) and
\(y_k=r_k/\lVert r_k\rVert\).

\uline{Guarantee}: At restart \(k\), \(x_{k+1}\) is the unique minimizer of
\(f\) over \(x_k+\mathcal K_2(A,r_k)\), and
\(r_{k+1}\perp\mathcal K_2(A,r_k)\).

\(y_0\gets r_0/\lVert r_0\rVert\).  For \(k=0,1,2,\ldots\), perform:
\begin{enumerate}[leftmargin=*]
  \item At restart \(k\), compute the named vectors \(Ar_k\) and
  \(A^2r_k\).
  \item Solve for \(\gamma_{k,0},\gamma_{k,1}\in\mathbb R\) so that
  \[
    x_{k+1}=x_k+\gamma_{k,0}r_k+\gamma_{k,1}Ar_k
  \]
  minimizes \(f\) on \(x_k+\mathcal K_2(A,r_k)\).
  \item Set
  \[
    r_{k+1}=r_k-\gamma_{k,0}Ar_k-\gamma_{k,1}A^2r_k,
    \qquad
    y_{k+1}=\frac{r_{k+1}}{\lVert r_{k+1}\rVert}.
  \]
\end{enumerate}
\end{algbox}

\begin{lemma}[Residual-polynomial form and implementation cost]
\label{lem:iteration-equivalence}
For every nonzero residual \(r_k\) of
\hyperref[alg:restarted-cg2]{\textsc{Restarted-CG(2)}} having grade at least
three, there is a unique degree-two polynomial \(P_k\) with \(P_k(0)=1\)
such that
\[
 r_{k+1}=P_k(A)r_k,\qquad
 r_{k+1}\perp\operatorname{span}\{r_k,Ar_k\}.
\]
Writing \(y_k=r_k/\lVert r_k\rVert\), one has
\(P_k=P_{y_k}\) and \(y_{k+1}=T(y_k)\).  If a multiplication by \(A\)
costs \(T_A\), then a restart uses \(2T_A+O(n)\) arithmetic operations;
for an explicit dense matrix the cost is \(O(n^2)\).
\end{lemma}

\begin{proof}
Since
\(x_{k+1}-x_k\in\operatorname{span}\{r_k,Ar_k\}\), the residual has the
form
\[
 r_{k+1}=r_k-A(x_{k+1}-x_k)
 =(I+a_kA+b_kA^2)r_k.
\]
First-order optimality for the strictly convex function \(f\) says exactly
that \(r_{k+1}\perp\operatorname{span}\{r_k,Ar_k\}\); conversely, these two
orthogonality equations characterize the minimizer.  Their moment matrix is
nonsingular when the grade is at least three, so the constant-normalized
polynomial \(P_k(t)=1+a_kt+b_kt^2\) is unique.  Also \(b_k\neq0\):
otherwise the two equations for \(1+a_kt\) would imply
\(\mu_0(y_k)\mu_2(y_k)=\mu_1(y_k)^2\), contrary to strict
Cauchy--Schwarz.  Hence \(P_k\) has degree two.

Multiplying \(P_k\) by a nonzero scalar gives the unique monic quadratic
orthogonal to \(1,t\) for \(\nu_{y_k}\).  It is therefore
\(\pi_{y_k}\).  Its constant term is positive, so
\(P_k=\pi_{y_k}/c_0(y_k)\).  Division by this positive scalar does not
change the normalized direction, which proves \(y_{k+1}=T(y_k)\) and the
algorithm's stated output guarantee.

Once \(Ar_k\) and \(A^2r_k\) have been computed, the two moment equations
for \(\gamma_{k,0},\gamma_{k,1}\), the constant-size solve, and the vector
updates take \(O(n)\) scalar operations.  This proves the resource bound.
\end{proof}

The polynomial representation permits a lossless reduction to the distinct
eigenvalues that occur in the initial residual.

\begin{lemma}[Reduction to distinct spectral coordinates]
\label{lem:spectral-reduction}
Let
\(\Lambda(A,r_0)=\{\lambda_1<\cdots<\lambda_m\}\), where
\(m=d(A,r_0)\).  There are
fixed orthonormal eigenvectors \(u_i\) and scalars \(z_{k,i}\) such that
\[
 Au_i=\lambda_i u_i,\qquad
 r_k=\sum_{i=1}^m z_{k,i}u_i
\]
for every \(k\).  Under the isometry
\(\sum_i z_i u_i\mapsto z\), the normalized residual iteration is exactly
\Cref{eq:monic-map} with
\(A=\operatorname{diag}(\lambda_1,\ldots,\lambda_m)\).  Thus convergence
in reduced coordinates is equivalent to convergence in the original space,
including when the original matrix has repeated eigenvalues.
\end{lemma}

\begin{proof}
Let \(\Pi_i=\Pi_{\lambda_i}\) and set
\(u_i=\Pi_ir_0/\lVert\Pi_ir_0\rVert\).  Every residual is a polynomial in
\(A\) applied to \(r_0\) by \Cref{lem:iteration-equivalence}; hence
\(\Pi_ir_k\) is always a scalar multiple of \(u_i\).  The displayed
expansion follows.  The stated map is an isometry intertwining \(A\) with
the diagonal matrix of distinct active eigenvalues, and so it preserves all
moment equations, projection polynomials, norms, and normalized directions.
\end{proof}

From now on, we work in these reduced coordinates.  Thus
\[
 A=\operatorname{diag}(\lambda_1,\ldots,\lambda_m),
 \qquad 0<\lambda_1<\cdots<\lambda_m.
\]
Coordinates that become zero remain zero.

\subsection{Two monotone quantities}

The first Lyapunov quantity is the norm of the monic quadratic projection.
Its exact two-step identity also controls the displacement of a fixed
parity of normalized residuals.

\begin{lemma}[Monic-projection Lyapunov identity]
\label{lem:monic-monotonicity}
Let \(y_{k+1}=T(y_k)\) and \(\theta_k=\theta(y_k)\).  As long as all
iterates have grade at least three,
\[
 \theta_k=\theta_{k+1}\langle y_k,y_{k+2}\rangle.
\tag{3}\label{eq:theta-identity}
\]
Consequently, \(\{\theta_k\}\) is nondecreasing and converges to a finite
\(\tau>0\), and
\[
 \lVert y_{k+2}-y_k\rVert^2
 =2\left(1-\frac{\theta_k}{\theta_{k+1}}\right)\longrightarrow0.
\tag{4}\label{eq:two-step-displacement}
\]
In fact,
\(\sum_{k\geq0}\lVert y_{k+2}-y_k\rVert^2<\infty\).
\end{lemma}

\begin{proof}
Put \(\pi_k=\pi_{y_k}\).  Since
\(\pi_k(A)y_k\perp\operatorname{span}\{y_k,Ay_k\}\), every monic
quadratic \(q\) satisfies
\[
 \begin{aligned}
 \theta_k^2
 &=\langle\pi_k(A)y_k,A^2y_k\rangle\\
 &=\langle\pi_k(A)y_k,q(A)y_k\rangle
 =\theta_k\langle y_{k+1},q(A)y_k\rangle.
 \end{aligned}
\]
Symmetry of \(A\) and division by \(\theta_k>0\) give
\(\theta_k=\langle y_k,q(A)y_{k+1}\rangle\).  Taking
\(q=\pi_{k+1}\) proves \Cref{eq:theta-identity}.
The inner product in that identity is positive and at most one, so
\(0<\theta_k\leq\theta_{k+1}\).  Taking \(t^2\) as a trial polynomial in
\Cref{eq:monic-minimization} yields \(\theta_k\leq\lVert A\rVert^2\);
hence \(\theta_k\to\tau>0\).

Both vectors in \Cref{eq:two-step-displacement} have unit norm, and
\Cref{eq:theta-identity} gives the formula.  Finally,
\(1-a/b\leq\log(b/a)\) for \(0<a\leq b\), so
\[
 \sum_{k=0}^N\lVert y_{k+2}-y_k\rVert^2
 \leq2\sum_{k=0}^N\log\frac{\theta_{k+1}}{\theta_k}
 =2\log\frac{\theta_{N+1}}{\theta_0}.
\]
The right side is uniformly bounded.
\end{proof}

The second monotone quantity uses the constant-normalized residual
polynomials.  Besides supplying another limiting scalar, it proves that the
iteration never terminates under the grade hypothesis.

\begin{lemma}[Energy-ratio monotonicity and grade persistence]
\label{lem:energy-ratio}
Assume \(d(A,r_0)\geq3\).  Every residual generated by
\hyperref[alg:restarted-cg2]{\textsc{Restarted-CG(2)}} is nonzero and has
grade at least three.  With
\[
 E_k=\lVert r_k\rVert_{A^{-1}}^2,\qquad
 \rho_k=\frac{E_{k+1}}{E_k},
\]
one has
\[
 0<\rho_k\leq\rho_{k+1}\leq
 \left(\frac{\lambda_{\max}-\lambda_{\min}}
 {\lambda_{\max}+\lambda_{\min}}\right)^2<1,
\tag{5}\label{eq:energy-ratio-bound}
\]
where
\(\lambda_{\min}=\min\Lambda(A,r_0)\) and
\(\lambda_{\max}=\max\Lambda(A,r_0)\).  In particular,
\(\rho_k\to\rho_\infty\) for some \(\rho_\infty\in(0,1)\).
\end{lemma}

\begin{proof}
For a nonzero residual, the minimizing step has a residual polynomial of
degree at most two with constant term one; if the residual has grade at most
two, one may choose such a polynomial to annihilate it.  Let \(P_k(0)=1\)
denote the minimizing polynomial.  Whenever \(r_{k+1}\neq0\),
orthogonality gives
\[
 \langle r_k,r_{k+2}\rangle_{A^{-1}}=E_{k+1}.
\tag{6}\label{eq:energy-cross-identity}
\]
Indeed,
\[
 \begin{aligned}
 \langle r_k,r_{k+2}\rangle_{A^{-1}}
 &=\langle P_{k+1}(A)r_k,r_{k+1}\rangle_{A^{-1}}\\
 &=E_{k+1}
 +\langle(P_{k+1}-P_k)(A)r_k,r_{k+1}\rangle_{A^{-1}}.
 \end{aligned}
\]
Since \(P_{k+1}(0)=P_k(0)=1\), their difference is \(t h(t)\) for a
polynomial \(h\) of degree at most one.  The final term is
\(\langle h(A)r_k,r_{k+1}\rangle=0\), proving
\Cref{eq:energy-cross-identity}.

If \(r_k\) has grade at least three, then \(r_{k+1}\neq0\), because a
degree-two polynomial cannot annihilate it.  If \(r_{k+1}\) had grade at
most two, the following minimizing step would give \(r_{k+2}=0\), whereas
\Cref{eq:energy-cross-identity} would give \(0=E_{k+1}>0\).
Induction proves nonvanishing and grade at least three for every residual.

Cauchy--Schwarz applied to \Cref{eq:energy-cross-identity} yields
\(E_{k+1}^2\leq E_kE_{k+2}\), which is
\(\rho_k\leq\rho_{k+1}\).  The energy decreases strictly at a nonzero
step: the admissible direction \(r_k\) has directional derivative
\(\langle\nabla f(x_k),r_k\rangle=-\lVert r_k\rVert^2<0\), so \(x_k\)
is not the affine-space minimizer.  Thus \(0<\rho_k<1\).  Since the CG
minimizer also minimizes the
\(A^{-1}\)-norm of the residual over degree-two residual polynomials with
constant term one, it does at least as well as
\[
 P(t)=1-\frac{2t}{\lambda_{\min}+\lambda_{\max}}.
\]
On the initial cyclic subspace,
\[
 |P(\lambda)|\leq
 \frac{\lambda_{\max}-\lambda_{\min}}
      {\lambda_{\max}+\lambda_{\min}}.
\]
Evaluation componentwise in the \(A^{-1}\)-norm proves
\Cref{eq:energy-ratio-bound}, and monotone convergence gives the claimed
limit.
\end{proof}

\subsection{Omega-limit structure and limiting quartics}

We next convert the two scalar limits into a rigid description of every
subsequential limit.

\begin{lemma}[Two-cycle structure of the omega-limit set]
\label{lem:omega-structure}
Let \(\Omega=\omega(y_{2k})\).
Then \(\Omega\) is nonempty, compact, and connected.  Every
\(y\in\Omega\) has grade three or four and satisfies
\[
 T^2y=y,\qquad \theta(y)=\theta(Ty)=\tau.
\tag{7}\label{eq:omega-two-cycle}
\]
If
\[
 q_y(t)=\pi_{Ty}(t)\pi_y(t),\qquad \beta=\tau^2,
\]
then
\[
 q_y(A)y=\beta y.
\tag{8}\label{eq:omega-quartic-eigen}
\]
\end{lemma}

\begin{proof}
The sequence lies on the compact unit sphere, and
\Cref{lem:monic-monotonicity} gives
\(\lVert y_{2k+2}-y_{2k}\rVert\to0\).  The limit-set facts recorded in
\Cref{sec:preliminaries} therefore show that \(\Omega\) is nonempty,
compact, and connected.

Take \(y_{2k_j}\to y\in\Omega\).  The limit
\(\theta(y_{2k_j})\to\tau>0\) rules out grade at most two.  Indeed, a fixed
monic quadratic \(g\) vanishing on such a support would satisfy
\[
 0<\theta(y_{2k_j})\leq\lVert g(A)y_{2k_j}\rVert
 \longrightarrow\lVert g(A)y\rVert=0,
\]
contradicting \Cref{eq:monic-minimization}.  Hence the moment system is
nonsingular at \(y\), and continuity gives
\[
 \theta(y)=\tau,\qquad y_{2k_j+1}\to Ty.
\]
The same argument at \(Ty\) gives grade at least three,
\(\theta(Ty)=\tau\), and continuity there.  Thus
\(y_{2k_j+2}\to T^2y\); together with
\(y_{2k_j+2}-y_{2k_j}\to0\), this proves
\Cref{eq:omega-two-cycle}.  Finally,
\[
 q_y(A)y
 =\pi_{Ty}(A)\bigl(\theta(y)Ty\bigr)
 =\theta(y)\theta(Ty)T^2y
 =\beta y.
\]
Every eigenvalue in the support of \(y\) is therefore a root of the
nonzero monic quartic \(q_y-\beta\).  The support has between three and
four points.
\end{proof}

The next lemma makes the two-step polynomial independent of the chosen
omega-limit point.  The energy-ratio limit fixes the otherwise undetermined
constant coefficient on a three-point support.

\begin{lemma}[Convergence of the two-step quartic]
\label{lem:quartic-convergence}
There is a monic quartic
\[
 q_\star(t)=t^4+\alpha_3^\star t^3+\alpha_2^\star t^2
             +\alpha_1^\star t+\alpha_0^\star
\]
such that the coefficients of
\[
 q_{2k}(t)=\pi_{y_{2k+1}}(t)\pi_{y_{2k}}(t)
\]
converge to those of \(q_\star\).  There are also monic quadratics
\(\pi_\star^{\mathrm e}\) and \(\pi_\star^{\mathrm o}\) such that
\[
 \pi_{y_{2k}}\longrightarrow\pi_\star^{\mathrm e},\qquad
 \pi_{y_{2k+1}}\longrightarrow\pi_\star^{\mathrm o},\qquad
 q_\star=\pi_\star^{\mathrm o}\pi_\star^{\mathrm e}
\]
coefficientwise.  Moreover,
\[
 \alpha_0^\star=\frac{\beta}{\rho_\infty}.
\tag{9}\label{eq:quartic-constant}
\]
\end{lemma}

\begin{proof}
For \(y\in\Omega\), write
\[
 q_y(t)=t^4+\alpha_3(y)t^3+\alpha_2(y)t^2+
             \alpha_1(y)t+\alpha_0(y).
\]
The coefficients are continuous in \(y\), and
\Cref{lem:omega-structure} gives \(q_y(\lambda_i)=\beta\) on
\(\operatorname{supp}(y)\).

Because \(P_y=\pi_y/\pi_y(0)\), the actual two-step residual polynomial at
the two-cycle \(y\) is \(q_y/\alpha_0(y)\).  By
\Cref{eq:omega-quartic-eigen}, it multiplies \(y\) by the positive scalar
\(\beta/\alpha_0(y)\).  Define the energy contraction functional
\[
 \mathcal E(z)=\frac{\lVert P_z(A)z\rVert_{A^{-1}}^2}
                    {\lVert z\rVert_{A^{-1}}^2}.
\]
Along a subsequence converging to \(y\), continuity and
\Cref{lem:energy-ratio} give
\(\mathcal E(y)=\mathcal E(Ty)=\rho_\infty\).  Hence
\[
 \left(\frac{\beta}{\alpha_0(y)}\right)^2=\rho_\infty^2.
\]
All quantities are positive, so
\(\alpha_0(y)=\beta/\rho_\infty\) throughout \(\Omega\).

For a fixed four-element support, the four equations
\(q_y(\lambda_i)=\beta\) uniquely determine the four lower coefficients of
the monic quartic.  For a fixed three-element support
\(\{i_1,i_2,i_3\}\), the constant coefficient is already fixed, and the
remaining system is nonsingular because
\[
 \det[\lambda_{i_a}^{\,b}]_{a,b=1}^3
 =\lambda_{i_1}\lambda_{i_2}\lambda_{i_3}
   \prod_{1\leq a<b\leq3}(\lambda_{i_b}-\lambda_{i_a})\neq0.
\]
Only finitely many three- and four-element supports occur.  Thus the
continuous coefficient map has finite image on the connected set
\(\Omega\), so it is constant there; call its value \(q_\star\).
Compactness then promotes subsequential convergence to convergence of the
full sequence \(q_{2k}\).

For every \(y\in\Omega\),
\(q_\star=\pi_{Ty}\pi_y\), so \(\pi_y\) is a monic quadratic divisor of
\(q_\star\).  A fixed quartic has finitely many monic divisors, whereas
\(y\mapsto\pi_y\) is continuous on connected \(\Omega\).  Therefore
\(\pi_y\) is one fixed polynomial \(\pi_\star^{\mathrm e}\) on
\(\Omega\), and compactness gives
\(\pi_{y_{2k}}\to\pi_\star^{\mathrm e}\).
Since \(q_{2k}=\pi_{y_{2k+1}}\pi_{y_{2k}}\) and all factors are monic,
coefficient comparison yields
\(\pi_{y_{2k+1}}\to
q_\star/\pi_\star^{\mathrm e}=:\pi_\star^{\mathrm o}\).
\end{proof}

The limiting quartic makes every coordinate outside its equality set
summably small, even if its limiting two-step multiplier is \(-1\).

\begin{lemma}[Summability away from the quartic equality set]
\label{lem:off-support-summability}
Let
\[
 I_\star=\{i:q_\star(\lambda_i)=\beta\}.
\]
For every \(i\notin I_\star\),
\[
 \sum_{k=0}^{\infty}y_{2k,i}^2<\infty.
\tag{10}\label{eq:off-support-summability}
\]
Consequently, the total squared mass outside \(I_\star\) is summable.
\end{lemma}

\begin{proof}
Put \(z_k=y_{2k}\).  Two applications of \Cref{eq:monic-map} give
\[
 z_{k+1,i}=s_{k,i}z_{k,i},\qquad
 s_{k,i}=
 \frac{q_{2k}(\lambda_i)}{\theta_{2k}\theta_{2k+1}}.
\]
By \Cref{lem:monic-monotonicity,lem:quartic-convergence},
\[
 s_{k,i}\longrightarrow
 s_{\star,i}=\frac{q_\star(\lambda_i)}{\beta}.
\]
If \(i\notin I_\star\), then \(s_{\star,i}\neq1\).  For some
\(\delta_i>0\) and all sufficiently large \(k\),
\(\lvert s_{k,i}-1\rvert\geq\delta_i\), and therefore
\[
 \lvert z_{k+1,i}-z_{k,i}\rvert^2
 \geq\delta_i^2\lvert z_{k,i}\rvert^2.
\]
The left side is summable by \Cref{lem:monic-monotonicity}.  This proves
\Cref{eq:off-support-summability}; summing over the finite spectrum gives
the final assertion.
\end{proof}

\subsection{Collapse on three spectral coordinates}

If all omega-limit points have grade three, the fixed quartic and the two
orthogonality equations leave only finitely many possible directions.

\begin{lemma}[Collapse of the three-coordinate regime]
\label{lem:three-point-collapse}
For
\[
 I_\star=\{i:q_\star(\lambda_i)=\beta\}
\]
one has \(3\leq\lvert I_\star\rvert\leq4\).  If
\(\lvert I_\star\rvert=3\), then \(\{y_{2k}\}\) converges.  More generally,
the same conclusion holds if every vector in \(\Omega\) has grade three.
In either case, there is a vector \(y_\star\) such that
\[
 y_{2k}\to y_\star,\qquad y_{2k+1}\to T(y_\star).
\]
\end{lemma}

\begin{proof}
Every \(y\in\Omega\) has at least three support points, all in
\(I_\star\), by
\Cref{lem:omega-structure,lem:quartic-convergence}.  Since
\(q_\star-\beta\) is a nonzero monic quartic,
\(\lvert I_\star\rvert\leq4\).  Thus
\(\lvert I_\star\rvert=3\) implies that every vector in \(\Omega\) has
grade three.

Assume more generally that all vectors in \(\Omega\) have grade three.
For each such vector,
\[
 q_\star=\pi_{Ty}\pi_y.
\]
Thus \(\pi_y\) is one of finitely many monic quadratic divisors of
\(q_\star\).  Fix a possible support
\(S=\{i_1,i_2,i_3\}\subseteq I_\star\) and a possible divisor \(\pi\).
For a vector \(y\in\Omega\) with this support and \(\pi_y=\pi\), put
\[
 w_a=y_{i_a}^2>0,\qquad v_a=w_a\pi(\lambda_{i_a}).
\]
Orthogonality and normalization become
\[
 \sum_{a=1}^3v_a=0,\qquad
 \sum_{a=1}^3\lambda_{i_a}v_a=0,\qquad
 \sum_{a=1}^3w_a=1.
\]
None of \(\pi(\lambda_{i_a})\) is zero, because its product with the other
factor equals \(\beta>0\).  The nullspace of the first two equations is
one-dimensional, spanned by
\[
 \bigl(\lambda_{i_3}-\lambda_{i_2},
       -(\lambda_{i_3}-\lambda_{i_1}),
       \lambda_{i_2}-\lambda_{i_1}\bigr).
\]
Hence the positive weights \(w_a=v_a/\pi(\lambda_{i_a})\) are unique after
normalization.  There are finitely many supports, divisors, and coordinate
sign choices, so \(\Omega\) is finite.  It is nonempty and connected by
\Cref{lem:omega-structure}, hence is a singleton \(\{y_\star\}\).
Precompactness then gives \(y_{2k}\to y_\star\), and continuity of \(T\)
gives the odd convergence.
\end{proof}

\subsection{Four-node barycentric algebra}

The remaining case has four equality nodes.  We first isolate the exact
one-step algebra on a four-node face; its barycentric parameter is the
coordinate later used to control drift along a continuum of two-cycles.
For a polynomial \(h\), the notation \([t^j]h(t)\) denotes the coefficient
of \(t^j\) in \(h\).

\begin{lemma}[Exact four-node barycentric transport]
\label{lem:four-node-coordinate}
Fix \(0<\xi_1<\xi_2<\xi_3<\xi_4\) and positive weights
\(w_1,\ldots,w_4\) summing to one.  Let \(\pi\) be the monic quadratic
satisfying
\[
 \sum_iw_i\pi(\xi_i)=0,\qquad
 \sum_i\xi_iw_i\pi(\xi_i)=0,
\]
put \(H=\sum_iw_i\pi(\xi_i)^2\), and define
\[
 w_i'=\frac{w_i\pi(\xi_i)^2}{H}.
\]
Let \(\pi'\) and \(H'\) be the analogous monic quadratic and squared
projection norm for \(w'\).  Set
\[
 L(t)=\prod_{i=1}^4(t-\xi_i),\qquad
 D_i=L'(\xi_i),\qquad q(t)=\pi(t)\pi'(t).
\]
There are unique real numbers \(\zeta,\zeta'\) such that
\[
 w_i\pi(\xi_i)=H\frac{\xi_i-\zeta}{D_i},\qquad
 w_i'\pi'(\xi_i)=H'\frac{\xi_i-\zeta'}{D_i}.
\tag{11}\label{eq:four-node-barycentric}
\]
If
\[
 \delta=[t^3]q(t)+\sum_{i=1}^4\xi_i,
\]
then
\[
 q(t)=L(t)+H'
 +\delta\,\frac{L(t)-L(\zeta)}{t-\zeta},
\qquad
 \zeta'-\zeta=\delta\,\frac{L(\zeta)}{H'}.
\tag{12}\label{eq:four-node-transport}
\]
In particular, \(\delta=0\) implies that two normalized polynomial steps
return every signed vector with squared coordinates \(w_i\) to itself.
\end{lemma}

\begin{proof}
The polynomial \(\pi\) cannot vanish at two of the four nodes.  Otherwise,
the two orthogonality equations on the remaining two distinct nodes force
both remaining quantities \(w_i\pi(\xi_i)\) to vanish, so \(\pi\) would
vanish at all four nodes.  Hence \(w'\) has at least three positive
coordinates, and \(\pi'\) and \(H'>0\) are well defined.

Lagrange interpolation gives, for every polynomial \(g\) of degree at most
three,
\[
 [t^3]g(t)=\sum_{i=1}^4\frac{g(\xi_i)}{D_i}.
\tag{13}\label{eq:lagrange-leading}
\]
The vector \(v_i=w_i\pi(\xi_i)\) annihilates the rows
\((1,\ldots,1)\) and \((\xi_1,\ldots,\xi_4)\).  By
\Cref{eq:lagrange-leading}, every vector in this two-dimensional nullspace
has a unique representation
\[
 v_i=\frac{a\xi_i+b}{D_i}.
\]
Multiplying by \(\pi(\xi_i)\) and summing, then applying
\Cref{eq:lagrange-leading} to \((at+b)\pi(t)\), gives
\(a=\sum_iw_i\pi(\xi_i)^2=H\).  Writing \(b=-H\zeta\) proves the first
identity in \Cref{eq:four-node-barycentric}; the second is identical.

The weight update and \Cref{eq:four-node-barycentric} imply
\[
 (\xi_i-\zeta)q(\xi_i)=H'(\xi_i-\zeta')
\]
at every node.  Hence the monic degree-five polynomial
\[
 (t-\zeta)q(t)-H'(t-\zeta')
\]
is divisible by \(L(t)\).  Comparison of its degree-four coefficient shows
that the quotient is \(t-\zeta+\delta\), so
\[
 (t-\zeta)q(t)-H'(t-\zeta')
 =(t-\zeta+\delta)L(t).
\tag{14}\label{eq:four-node-divisibility}
\]
Putting \(t=\zeta\) gives the second identity in
\Cref{eq:four-node-transport}; substituting it into
\Cref{eq:four-node-divisibility} and dividing by \(t-\zeta\) gives the
first.

If \(\delta=0\), then \(q=L+H'\), and \(q(\xi_i)=H'>0\).  After two
weight updates,
\[
 w_i''=\frac{w_iq(\xi_i)^2}{HH'}=w_i\frac{H'}{H}.
\]
Both \(w\) and \(w''\) sum to one, so \(H'=H\) and \(w''=w\).  The signed
two-step multiplier is \(q(\xi_i)/\sqrt{HH'}=1\), proving the last
assertion.
\end{proof}

Assume from now on that \(\lvert I_\star\rvert=4\).
\Cref{lem:monic-monotonicity,lem:quartic-convergence} give
\(\theta_n\to\tau>0\), convergence of the even and odd monic factors, and
the limiting quartic \(q_\star\), while
\Cref{lem:off-support-summability} gives
\[
 \sum_k\sum_{j\notin I_\star}y_{2k,j}^2<\infty.
\]
We retain the contribution of every coordinate outside the four-node
equality set.

\begin{definition}[Four-point orbit data]
\label{def:four-point-setting}
A \emph{four-point orbit} is the reduced diagonal iteration
\(y_{n+1}=T(y_n)\) in the case \(\lvert I_\star\rvert=4\).
For such an orbit, write the four indices as
\(\iota_1,\ldots,\iota_4\), ordered so that
\[
 \xi_1:=\lambda_{\iota_1}<\xi_2:=\lambda_{\iota_2}
 <\xi_3:=\lambda_{\iota_3}<\xi_4:=\lambda_{\iota_4}.
\]
For \(n\geq0\), set
\[
 p_n=\pi_{y_n},\qquad \theta_n=\theta(y_n),\qquad
 H_n=\theta_n^2,\qquad q_n=p_np_{n+1},
\]
\[
 v_{n,i}=y_{n,\iota_i}^2\quad(1\leq i\leq4),\qquad
 w_{n,j}=y_{n,j}^2\quad(j\notin I_\star).
\]
Let
\[
 p_{\mathrm e}=\pi_\star^{\mathrm e},\qquad
 p_{\mathrm o}=\pi_\star^{\mathrm o},\qquad
 L(t)=\prod_{i=1}^4(t-\xi_i),\qquad D_i=L'(\xi_i),
\]
and put
\[
 \beta=\tau^2,\qquad
 \delta_n=[t^3]q_n+\sum_{i=1}^4\xi_i,\qquad
 \epsilon_n=\sum_{j\notin I_\star}w_{n,j}.
\]
\end{definition}

The quadratic factors converge separately by
\Cref{lem:quartic-convergence}; on odd indices they are merely interchanged.
Thus \(q_n\to q_\star\) along the full sequence.  Since
\(q_\star-\beta\) and \(L\) are monic quartics with the same four roots,
\[
 q_\star(t)=L(t)+\beta=p_{\mathrm o}(t)p_{\mathrm e}(t).
\tag{15}\label{eq:four-point-limit-quartic}
\]
In particular, \(H_n\to\beta\) and \(\delta_n\to0\).

\begin{definition}[Tail and resonance data]
\label{def:four-point-tail-data}
Define
\[
 K_n=\sqrt{H_nH_{n+1}}=\theta_n\theta_{n+1},\qquad
 \Delta_n=\lVert y_{n+2}-y_n\rVert^2,
\]
\[
 \sigma_n=\frac{H_{n+1}}{K_n}-1,\qquad
 e_n=\frac{\beta}{K_n}-1,\qquad
 b_{n,i}=\frac{q_n(\xi_i)}{K_n}-1.
\]
Choose \(N_0\) beyond the first zero time of every coordinate that is ever
annihilated.  Among indices \(j\notin I_\star\) for which \(w_{n,j}>0\)
for all \(n\geq N_0\), define
\[
 \mathcal R=\{j:q_\star(\lambda_j)=-\beta\},\qquad
 \mathcal S=\{j:\lvert q_\star(\lambda_j)\rvert<\beta\},
\]
and put
\[
 \epsilon_n^{\mathrm R}=\sum_{j\in\mathcal R}w_{n,j},\qquad
 \epsilon_n^{\mathrm S}=\sum_{j\in\mathcal S}w_{n,j}.
\]
\end{definition}

The following division functional extracts the scalar coefficient transported
from one quadratic factor to the next.

\begin{definition}[Division functional]
\label{def:division-functional}
For \(z\in\mathbb R\), set
\[
 R_z(t)=\frac{L(t)-L(z)}{t-z}.
\]
For a monic quadratic \(p\), let \(X(p,z)\) be the coefficient of \(t\)
in the remainder of \(R_z(t)\) upon division by \(p(t)\).
\end{definition}

If
\[
 L(t)=t^4-St^3+e_2t^2-e_3t+e_4,\qquad
 p(t)=t^2+at+b,
\]
then direct division gives
\[
 X(p,z)=z^2-(a+S)z+a^2+aS-b+e_2.
\tag{16}\label{eq:division-functional-formula}
\]

Before deriving the forced transport law, we transfer the even-index
summability from \Cref{lem:off-support-summability} to both parities.

\begin{lemma}[Full-parity outside-mass summability]
\label{lem:outside-mass-all-parities}
Every four-point orbit satisfies
\[
 \sum_{n=0}^{\infty}\epsilon_n<\infty.
\tag{17}\label{eq:outside-mass-summable}
\]
\end{lemma}

\begin{proof}
By \Cref{lem:off-support-summability},
\(\sum_k\epsilon_{2k}<\infty\).  The coefficients of \(p_{2k}\) are
bounded, the spectrum is finite, and \(H_{2k}\) is bounded away from zero.
The exact one-step squared-weight update therefore gives, for all
sufficiently large \(k\),
\[
 \epsilon_{2k+1}
 =\sum_{j\notin I_\star}
   \frac{w_{2k,j}p_{2k}(\lambda_j)^2}{H_{2k}}
 \leq C\epsilon_{2k}.
\]
Adding the finitely many initial terms proves the claim.
\end{proof}

The outside coordinates cannot simply be discarded: a multiplier tending
to \(-1\) need not decay geometrically.  Lagrange correction incorporates
all of them into an exact barycentric identity.

\begin{definition}[Corrected barycentric coordinate]
\label{def:corrected-barycentric}
For \(x\notin\{\xi_1,\ldots,\xi_4\}\), define
\[
 \ell_i(x)=\frac{L(x)}{D_i(x-\xi_i)}
 \qquad(1\leq i\leq4).
\]
Put
\[
 a_{n,i}=v_{n,i}p_n(\xi_i)
 +\sum_{j\notin I_\star}
       w_{n,j}p_n(\lambda_j)\ell_i(\lambda_j).
\]
The corrected barycentric coordinate \(\zeta_n\) is the real number
satisfying
\[
 a_{n,i}=H_n\frac{\xi_i-\zeta_n}{D_i}
 \qquad(1\leq i\leq4).
\tag{18}\label{eq:corrected-barycentric}
\]
\end{definition}

\begin{lemma}[Forced barycentric transport]
\label{lem:forced-barycentric}
The coordinate \(\zeta_n\) in
\Cref{def:corrected-barycentric} exists and is unique.  Define
\[
 F_n(t)=\sum_{j\notin I_\star}
 w_{n,j}p_n(\lambda_j)L(\lambda_j)
 \frac{q_n(\lambda_j)-q_n(t)}{\lambda_j-t},
\]
where the divided difference is interpreted by polynomial continuation at
\(t=\lambda_j\).  Then
\begin{align}
 &(t-\zeta_n)q_n(t)+\frac{F_n(t)}{H_n}
       -H_{n+1}(t-\zeta_{n+1})
   =(t-\zeta_n+\delta_n)L(t),
 \tag{19}
 \label{eq:full-polynomial-identity}\\
 &\zeta_{n+1}-\zeta_n
   =\frac{\delta_nL(\zeta_n)-F_n(\zeta_n)/H_n}{H_{n+1}},
 \tag{20}
 \label{eq:full-zeta-increment}\\
 &q_n(t)=L(t)+H_{n+1}
  +\delta_n\frac{L(t)-L(\zeta_n)}{t-\zeta_n}
  -\frac{F_n(t)-F_n(\zeta_n)}
         {H_n(t-\zeta_n)}.
 \tag{21}
 \label{eq:full-quartic-identity}
\end{align}
Let
\[
 C_n(t)=\frac{F_n(t)-F_n(\zeta_n)}
              {H_n(t-\zeta_n)},\qquad
 Y_n(p)=[t]\operatorname{rem}_{p}C_n(t),
\]
again using polynomial continuation.  The exact defect transport law is
\[
 \delta_{n+1}X(p_{n+1},\zeta_{n+1})-Y_{n+1}(p_{n+1})
 =
 \delta_nX(p_{n+1},\zeta_n)-Y_n(p_{n+1}).
\tag{22}\label{eq:full-defect-transport}
\]
Uniformly for all sufficiently large \(n\),
\[
 \lVert F_n\rVert_{\mathrm{coeff}}+
 \lVert C_n\rVert_{\mathrm{coeff}}+
 \lvert Y_n(p_{n+1})\rvert\leq C\epsilon_n,
\tag{23}\label{eq:forcing-coefficient-bound}
\]
where \(\lVert P\rVert_{\mathrm{coeff}}\) is the largest absolute
coefficient of \(P\).  Finally, if
\(p_n(t)=t^2+c_{1,n}t+c_{0,n}\), then
\[
 \zeta_n=\sum_{i=1}^4\xi_i+c_{1,n}
          -\sum_{j=1}^m\lambda_jy_{n+1,j}^2.
\tag{24}\label{eq:zeta-state-formula}
\]
\end{lemma}

\begin{proof}
Lagrange interpolation gives
\[
 \sum_{i=1}^4\xi_i^r\ell_i(x)=x^r
 \qquad(0\leq r\leq3).
\tag{25}\label{eq:lagrange-moments}
\]
The two orthogonality equations for \(p_n\) imply
\[
 \sum_i a_{n,i}=0,\qquad \sum_i\xi_i a_{n,i}=0.
\]
Every four-vector satisfying these equations has the form
\((\kappa\xi_i+\nu)/D_i\).  Since \(t^2-p_n(t)\) is linear and \(p_n\)
is orthogonal to every linear polynomial,
\[
 \sum_i\xi_i^2a_{n,i}
 =\sum_{j=1}^m\lambda_j^2y_{n,j}^2p_n(\lambda_j)=H_n.
\]
The identities
\(\sum_i\xi_i^2/D_i=0\) and
\(\sum_i\xi_i^3/D_i=1\) show that \(\kappa=H_n\), proving existence and
uniqueness of \(\zeta_n\).

The exact squared-weight update is
\[
 y_{n+1,j}^2
 =\frac{y_{n,j}^2p_n(\lambda_j)^2}{H_n}.
\tag{26}\label{eq:squared-weight-update}
\]
Use it in the corrected formula at time \(n+1\) and substitute the formula
at time \(n\).  At each active node \(\xi_i\), this gives
\[
 H_{n+1}\frac{\xi_i-\zeta_{n+1}}{D_i}
 =
 q_n(\xi_i)\frac{\xi_i-\zeta_n}{D_i}
 +\sum_{j\notin I_\star}
 \frac{w_{n,j}p_n(\lambda_j)}{H_n}
 \bigl(q_n(\lambda_j)-q_n(\xi_i)\bigr)\ell_i(\lambda_j).
\]
Consequently, the left side of
\Cref{eq:full-polynomial-identity} vanishes at all four roots of \(L\).
It is a monic polynomial of degree five, and comparison of the \(t^4\)
coefficient shows that its quotient by \(L\) is
\(t-\zeta_n+\delta_n\).  This proves
\Cref{eq:full-polynomial-identity}.  Evaluation at \(t=\zeta_n\) gives
\Cref{eq:full-zeta-increment}; substitution back gives
\Cref{eq:full-quartic-identity}.

Reduce \Cref{eq:full-quartic-identity} modulo \(p_{n+1}\) and take the
coefficient of \(t\) in the remainder.  Both
\(q_n=p_np_{n+1}\) and \(q_{n+1}=p_{n+1}p_{n+2}\) are divisible by
\(p_{n+1}\).  Comparing the identities at \(n\) and \(n+1\) proves
\Cref{eq:full-defect-transport}.

The spectrum is fixed, the coefficients of \(p_n,q_n\) are bounded, and
\(H_n\) is bounded away from zero.  Formula
\Cref{eq:zeta-state-formula}, proved below, shows that \(\zeta_n\) is
bounded.  Thus every coefficient of \(F_n,C_n\), and the indicated
remainder is at most \(C\epsilon_n\), proving
\Cref{eq:forcing-coefficient-bound}.

To prove \Cref{eq:zeta-state-formula}, take the third moment of
\Cref{eq:corrected-barycentric}:
\[
 H_n\left(\sum_i\xi_i-\zeta_n\right)
 =\sum_{j=1}^m\lambda_j^3y_{n,j}^2p_n(\lambda_j).
\]
Orthogonality and the displayed form of \(p_n\) give
\[
 \sum_{j=1}^m\lambda_jy_{n,j}^2p_n(\lambda_j)^2
 =
 \sum_{j=1}^m\lambda_j^3y_{n,j}^2p_n(\lambda_j)+c_{1,n}H_n.
\]
After division by \(H_n\), the left side is the spectral mean of
\(y_{n+1}\) by \Cref{eq:squared-weight-update}, which proves
\Cref{eq:zeta-state-formula}.
\end{proof}

\subsection{Lyapunov tails in the four-coordinate regime}

It remains to collapse \(\Omega\) when its positive equality set has four
nodes.  The four-point notation and the exact forced transport identities
of \Cref{def:four-point-setting,lem:forced-barycentric} now apply, because
\Cref{lem:monic-monotonicity,lem:quartic-convergence,%
lem:off-support-summability} supply their conditional hypotheses.

The scalar defect transport is parabolic.  The following elementary
dichotomy will later turn geometric forcing into convergence of the
barycentric coordinate.

\begin{lemma}[Parabolic sign-or-summability dichotomy]
\label{lem:parabolic-dichotomy}
Let \(u_n\to0\), and suppose that for constants \(C>0\) and
\(0<\gamma<1\),
\[
 \lvert u_{n+1}-u_n\rvert\leq C(u_n^2+\gamma^n)
\]
for every sufficiently large \(n\).  Then either \(u_n\) is eventually
one-signed, or \(\sum_n\lvert u_n\rvert<\infty\).
\end{lemma}

\begin{proof}
Choose \(\alpha\) with \(\sqrt\gamma<\alpha<1\).  For all sufficiently
large \(n\), once \(\lvert u_n\rvert\geq\alpha^n\), its sign never changes.
Indeed, as long as \(\lvert u_k\rvert\geq\alpha^k\), one has
\(\gamma^k=o(u_k^2)\), and smallness gives
\[
 \lvert u_{k+1}-u_k\rvert\leq C_1u_k^2<\frac12\lvert u_k\rvert,
 \qquad
 \frac1{\lvert u_{k+1}\rvert}
 \leq\frac1{\lvert u_k\rvert}+2C_1.
\]
If \(m>n\) were the first later index with
\(\lvert u_m\rvert<\alpha^m\), iteration would yield
\[
 \lvert u_m\rvert
 \geq\frac1{\alpha^{-n}+2C_1(m-n)}>\alpha^m
\]
for all sufficiently large \(n\).  Writing \(\ell=m-n\), the final
inequality is equivalent to
\[
 \alpha^\ell+2C_1\ell\alpha^{n+\ell}<1,
\]
uniformly for \(\ell\geq1\), a contradiction.  Therefore infinitely many
sign changes force \(\lvert u_n\rvert<\alpha^n\) eventually, and the latter
sequence is summable.
\end{proof}

\begin{definition}[Four-positive points and surviving coordinates]
\label{def:four-positive-surviving}
An omega-limit point is \emph{four-positive} if its four squared
coordinates on \(I_\star\) are positive.  A coordinate is
\emph{surviving} if it is nonzero at arbitrarily large times.  Since a zero
coordinate remains zero, every surviving coordinate is nonzero at all
times after the index \(N_0\) in \Cref{def:four-point-tail-data}.
\end{definition}

A resonant outside coordinate has limiting two-step multiplier \(-1\).
Its sign change makes its mass visible in the exact Lyapunov increase.

\begin{lemma}[Active sign geometry and resonant-tail coercivity]
\label{lem:resonant-geometry-tail}
Assume that \(\Omega\) contains a four-positive vector.  Then
\[
 \operatorname{sgn}p_{\mathrm e}(\xi_i)
 =\operatorname{sgn}p_{\mathrm o}(\xi_i)
 =(+,-,-,+)_i.
\tag{27}\label{eq:active-sign-pattern}
\]
Every cluster value of \(\zeta_n\) belongs to
\([\xi_2,\xi_3]\), and every resonant node \(\lambda_j\),
\(j\in\mathcal R\), belongs to
\((\xi_1,\xi_2)\cup(\xi_3,\xi_4)\).  Moreover, for all sufficiently large
\(n\),
\[
 \tau-\theta_n
 =\frac12\sum_{\ell=n}^{\infty}\theta_{\ell+1}\Delta_\ell
 \geq c\sum_{\ell=n}^{\infty}\epsilon_\ell^{\mathrm R}
\tag{28}\label{eq:resonant-tail-coercivity}
\]
for a constant \(c>0\).
\end{lemma}

\begin{proof}
The four-positive point and its one-step image make
\(p_{\mathrm e}\) and \(p_{\mathrm o}\) orthogonal quadratics for positive
weights on all four active nodes.  Each has two simple real roots in
\((\xi_1,\xi_4)\).  To see this, if such a quadratic had fewer than two
sign-changing roots in that interval, let \(g\) be the product of the
factors \(t-r\) over its sign-changing roots.  Then \(\deg g\leq1\), while
the product of the quadratic with \(g\) has one nonzero sign on all four
nodes, contradicting orthogonality to every linear polynomial.

By \Cref{eq:four-point-limit-quartic}, none of these roots lies outside
\([\xi_1,\xi_4]\) or in the middle gap, where \(L+\beta>0\).  If one
factor placed both roots in the same outer gap, it would be positive at all
four active nodes, contradicting its orthogonality to \(1\).  Each factor
therefore has one root in each outer gap and has the sign pattern
\Cref{eq:active-sign-pattern}.

Formula \Cref{eq:zeta-state-formula} makes \(\zeta_n\) bounded.  Along a
convergent subsequence, \(\epsilon_n\to0\) by
\Cref{lem:outside-mass-all-parities}, so the correction in
\Cref{eq:corrected-barycentric} vanishes.  Passing to the limit and using
\Cref{eq:active-sign-pattern} and nonnegativity of the active weights
forces every cluster value of \(\zeta_n\) into \([\xi_2,\xi_3]\).
For \(j\in\mathcal R\),
\[
 L(\lambda_j)=q_\star(\lambda_j)-\beta=-2\beta<0,
\]
and the sign of the four-root polynomial \(L\) places \(\lambda_j\) in an
outer gap.

Orthogonality of \(p_n\) to every linear polynomial gives
\[
 \sum_{j=1}^m y_{n,j}^2p_n(\lambda_j)p_{n+1}(\lambda_j)
 =\sum_{j=1}^m y_{n,j}^2p_n(\lambda_j)^2=H_n.
\]
Consequently,
\[
 \langle y_n,y_{n+2}\rangle=\frac{\theta_n}{\theta_{n+1}},
 \qquad
 \theta_{n+1}-\theta_n
 =\frac{\theta_{n+1}}2\Delta_n.
\tag{29}\label{eq:exact-theta-increment}
\]
Summation from \(n\) to infinity proves the equality in
\Cref{eq:resonant-tail-coercivity}.  For \(j\in\mathcal R\),
\[
 \frac{y_{\ell+2,j}}{y_{\ell,j}}
 =\frac{q_\ell(\lambda_j)}{K_\ell}\longrightarrow-1.
\]
The finite resonant set thus satisfies
\[
 \lvert y_{\ell+2,j}-y_{\ell,j}\rvert^2\geq w_{\ell,j}
\]
for all sufficiently large \(\ell\).  Since
\(\theta_{\ell+1}\to\tau>0\), substitution into the exact tail identity
proves the lower bound.
\end{proof}

The resonant mass is now charged to the Lyapunov tail, while strict-gap
coordinates decay geometrically.  Together these estimates control every
outside coordinate.

\begin{lemma}[Lyapunov control of displacement and outside-mass tails]
\label{lem:outside-tail}
If \(\Omega\) contains a four-positive vector, then there are constants
\(C>0\) and \(0<\gamma<1\) such that, for all sufficiently large \(n\),
\[
 \sum_{\ell=n}^{\infty}\Delta_\ell+
 \sum_{\ell=n}^{\infty}\epsilon_\ell
 \leq C\bigl(\tau-\theta_n+\gamma^n\bigr).
\tag{30}\label{eq:combined-tail}
\]
Moreover,
\[
 \sigma_n=\frac{\Delta_n}{2-\Delta_n}\leq C\Delta_n,\qquad
 e_n\geq\frac{\tau-\theta_n}{\tau}.
\tag{31}\label{eq:sigma-e-relations}
\]
\end{lemma}

\begin{proof}
Equation \Cref{eq:exact-theta-increment} yields
\[
 \tau-\theta_n
 =\frac12\sum_{\ell=n}^{\infty}\theta_{\ell+1}\Delta_\ell.
\tag{32}\label{eq:exact-displacement-tail}
\]
Because \(\theta_\ell\to\tau>0\),
\[
 \sum_{\ell=n}^{\infty}\Delta_\ell
 \leq C(\tau-\theta_n).
\tag{33}\label{eq:displacement-tail}
\]
\Cref{lem:resonant-geometry-tail} likewise gives
\[
 \sum_{\ell=n}^{\infty}\epsilon_\ell^{\mathrm R}
 \leq C(\tau-\theta_n).
\tag{34}\label{eq:resonant-mass-tail}
\]

For every surviving outside coordinate, the exact recurrence is
\[
 w_{n+2,j}
 =w_{n,j}\frac{q_n(\lambda_j)^2}{H_nH_{n+1}}.
\tag{35}\label{eq:outside-two-step}
\]
By \Cref{lem:off-support-summability}, \(w_{2k,j}\to0\).  If
\(\lvert q_\star(\lambda_j)\rvert>\beta\), the multiplier in
\Cref{eq:outside-two-step} is eventually greater than one on each parity,
contradicting this convergence on the even parity and its one-step
consequence on the odd parity.  Equality
\(q_\star(\lambda_j)=\beta\) would put \(j\) in \(I_\star\).  Hence every
surviving outside coordinate lies in \(\mathcal R\cup\mathcal S\).

For \(j\in\mathcal S\), the multiplier in
\Cref{eq:outside-two-step} is eventually bounded by a number strictly
smaller than one.  Applying the recurrence separately to the even and odd
parities and using finiteness of the spectrum gives one
\(0<\gamma<1\) such that
\[
 \epsilon_n^{\mathrm S}\leq C\gamma^n,\qquad
 \sum_{\ell=n}^{\infty}\epsilon_\ell^{\mathrm S}
 \leq C\gamma^n.
\tag{36}\label{eq:strict-mass-tail}
\]
Nonsurviving coordinates are zero after \(N_0\).
Combining \Cref{eq:displacement-tail}--\Cref{eq:strict-mass-tail}
proves \Cref{eq:combined-tail}.

Finally, \Cref{eq:exact-theta-increment} gives
\[
 \sigma_n=\frac{\theta_{n+1}}{\theta_n}-1
 =\frac{\Delta_n}{2-\Delta_n}.
\]
Also \(\theta_{n+1}\leq\tau\), so
\[
 e_n
 =\frac{\tau^2-\theta_n\theta_{n+1}}
        {\theta_n\theta_{n+1}}
 \geq\frac{\tau-\theta_n}{\tau}.
\]
This proves \Cref{eq:sigma-e-relations}.
\end{proof}

The next estimate converts displacement and outside mass into summable
variation of the cubic quartic defect.  The factor
\(\lvert L(\zeta_n)\rvert\) is essential near a grade-three endpoint, where
one active weight may vanish.

\begin{definition}[Endpoint defect objects]
\label{def:forced-objects}
For \(z\in\mathbb R\), put
\[
 B_i(z)=R_z(\xi_i)=\frac{L(z)}{z-\xi_i},
\]
with \(B_i(\xi_i)=D_i\) by continuity.  Retain \(F_n,C_n\), and
\(Y_n\) from \Cref{lem:forced-barycentric}.
\end{definition}

\begin{lemma}[Summable variation of the forced quartic defect]
\label{lem:defect-tail}
Assume that \(\Omega\) contains a four-positive vector.  For all
sufficiently large \(n\),
\[
 q_n(t)=L(t)+H_{n+1}+\delta_nR_{\zeta_n}(t)-C_n(t)
\tag{37}\label{eq:forced-quartic}
\]
and
\[
 \delta_{n+1}X(p_{n+1},\zeta_{n+1})-Y_{n+1}(p_{n+1})
 =
 \delta_nX(p_{n+1},\zeta_n)-Y_n(p_{n+1}).
\tag{38}\label{eq:forced-transport-local}
\]
There is a constant \(C>0\) such that
\begin{align}
 \delta_n^2\lvert L(\zeta_n)\rvert
 &\leq C(\Delta_n+\epsilon_n),
 \tag{39}
 \label{eq:defect-energy}\\
 \lvert\delta_{n+1}-\delta_n\rvert
 &\leq C(\Delta_n+\epsilon_n+\epsilon_{n+1}),
 \tag{40}
 \label{eq:defect-variation}\\
 \lvert\delta_n\rvert
 &\leq C\bigl(\tau-\theta_n+\gamma^n\bigr),
 \tag{41}
 \label{eq:defect-tail}
\end{align}
where \(\gamma\) is as in \Cref{lem:outside-tail}.
\end{lemma}

\begin{proof}
Equations \Cref{eq:forced-quartic} and
\Cref{eq:forced-transport-local} are
\Cref{eq:full-quartic-identity} and
\Cref{eq:full-defect-transport}.

We first record uniform nonvanishing of \(X\).  Let \(p\) be either
limiting quadratic factor and let \(\widehat p\) be the other, so that
\(p\widehat p=L+\beta\).  If \(u_1,u_2\) are the roots of \(p\), then the
coefficient of \(t\) in the linear interpolation of \(R_z\) at these roots
is
\[
 \frac{R_z(u_2)-R_z(u_1)}{u_2-u_1}
 =\frac{L(z)+\beta}{p(z)}
 =\widehat p(z).
\tag{42}\label{eq:X-limit}
\]
Thus \(X(p,z)=\widehat p(z)\).  Both roots of \(\widehat p\) lie in the
outer gaps by \Cref{lem:resonant-geometry-tail}, so \(\widehat p\) has no
zero on a fixed neighborhood of \([\xi_2,\xi_3]\).  Every cluster value of
\(\zeta_n\) lies in that interval, and the parity factors converge.
Consequently, for all sufficiently large \(n\), each of
\[
 \lvert X(p_{n+1},\zeta_n)\rvert,\qquad
 \lvert X(p_{n+1},\zeta_{n+1})\rvert
\]
is bounded below by a fixed positive constant.  The same functions are
uniformly Lipschitz in their second argument there.

Rearrange \Cref{eq:forced-transport-local}, use
\Cref{eq:forcing-coefficient-bound} at \(n\) and \(n+1\), and use
\Cref{eq:full-zeta-increment}.  This gives
\[
 \lvert\zeta_{n+1}-\zeta_n\rvert
 \leq C\bigl(\lvert\delta_n\rvert\lvert L(\zeta_n)\rvert+
              \epsilon_n\bigr)
\tag{43}\label{eq:zeta-forced-bound}
\]
and then
\[
 \lvert\delta_{n+1}-\delta_n\rvert
 \leq C\delta_n^2\lvert L(\zeta_n)\rvert+
       C(\epsilon_n+\epsilon_{n+1}).
\tag{44}\label{eq:variation-before-energy}
\]

It remains to prove \Cref{eq:defect-energy}.  Evaluation of
\Cref{eq:forced-quartic} at the active nodes gives
\[
 b_{n,i}
 =\sigma_n+\frac{\delta_n}{K_n}B_i(\zeta_n)
       -\frac{C_n(\xi_i)}{K_n}.
\tag{45}\label{eq:active-error-representation}
\]
Cover a fixed neighborhood of \([\xi_2,\xi_3]\) by small endpoint
neighborhoods and a compact interior region.  On the compact interior,
\(B_1(z)=L(z)/(z-\xi_1)\) is bounded away from zero directly, since
\(L\) is positive in the open middle gap.  Also \(v_{n,1}\) is uniformly
positive there.  Otherwise, a subsequence with
\(v_{n,1}\to0\) and \(\zeta_n\to z\) in that compact interior would make
the vanishing outside correction in \Cref{eq:corrected-barycentric} give
\[
 0=\beta\frac{\xi_1-z}{D_1},
\]
which is impossible.  Since
\[
 \Delta_n\geq v_{n,1}b_{n,1}^2,
\tag{46}\label{eq:active-energy}
\]
solving \Cref{eq:active-error-representation} at \(i=1\), using
\(\sigma_n=O(\Delta_n)\) and \(C_n=O(\epsilon_n)\), yields
\(\delta_n^2\leq C(\Delta_n+\epsilon_n)\) on the interior.  Boundedness of
\(L(\zeta_n)\) proves \Cref{eq:defect-energy} there.

Now suppose that \(\zeta_n\) is near an endpoint \(\xi_a\), where
\(a\in\{2,3\}\), and put
\[
 \omega_n=v_{n,a},\qquad d_n=b_{n,a}.
\]
The corrected barycentric equation at \(i=a\) is
\[
 H_n\frac{\xi_a-\zeta_n}{D_a}
 =\omega_np_n(\xi_a)
 +\sum_{j\notin I_\star}
   w_{n,j}p_n(\lambda_j)\ell_a(\lambda_j).
\tag{47}\label{eq:endpoint-barycentric}
\]
All polynomial values are bounded and \(H_n\) is bounded away from zero,
so
\[
 \lvert\zeta_n-\xi_a\rvert\leq C(\omega_n+\epsilon_n),
 \qquad
 \lvert L(\zeta_n)\rvert\leq C(\omega_n+\epsilon_n).
\tag{48}\label{eq:endpoint-distance}
\]
Moreover, \(B_a(\zeta_n)\) is bounded away from zero because
\(B_a(\xi_a)=D_a\neq0\).  Solving
\Cref{eq:active-error-representation} at \(i=a\) gives
\[
 \lvert\delta_n\rvert
 \leq C(\lvert d_n\rvert+\sigma_n+\epsilon_n).
\tag{49}\label{eq:delta-from-d}
\]
Here \(d_n\to0\), and
\[
 \omega_nd_n^2
 =\lvert y_{n+2,\iota_a}-y_{n,\iota_a}\rvert^2
 \leq\Delta_n.
\tag{50}\label{eq:endpoint-energy}
\]
Using \Cref{eq:endpoint-distance}--\Cref{eq:endpoint-energy},
\(\sigma_n=O(\Delta_n)\), and smallness of these quantities,
\[
 \begin{aligned}
 \delta_n^2\lvert L(\zeta_n)\rvert
 &\leq C(d_n^2+\sigma_n^2+\epsilon_n^2)
          (\omega_n+\epsilon_n)\\
 &\leq C\omega_nd_n^2+C\sigma_n^2+C\epsilon_n\\
 &\leq C(\Delta_n+\epsilon_n).
 \end{aligned}
\]
This proves \Cref{eq:defect-energy} near both endpoints and hence
everywhere for large \(n\).

Substitution into \Cref{eq:variation-before-energy} proves
\Cref{eq:defect-variation}.  We know \(\delta_n\to0\) from
\Cref{eq:four-point-limit-quartic}, while
\(\sum_n(\Delta_n+\epsilon_n)<\infty\) by
\Cref{lem:outside-tail}.  Therefore the variation estimate may be summed
backward from infinity:
\[
 \lvert\delta_n\rvert
 \leq C\sum_{\ell=n}^{\infty}(\Delta_\ell+\epsilon_\ell).
\]
\Cref{lem:outside-tail} now proves \Cref{eq:defect-tail}.
\end{proof}

The defect-tail bound controls the exceptional active multiplier during
visits near either endpoint of the middle gap.  All other active
multipliers are square-summable on those visits.

\begin{lemma}[Endpoint multiplier absorption]
\label{lem:endpoint-errors}
Assume that \(\Omega\) contains a four-positive vector, and fix
\(\gamma\) from \Cref{lem:outside-tail}.  There are
\(0<\varepsilon_0<(\xi_3-\xi_2)/3\), \(N\), and \(C>0\) such that, for
\(a\in\{2,3\}\), the endpoint-visit set
\[
 \mathcal N_a
 =\{n\geq N:\lvert\zeta_n-\xi_a\rvert<\varepsilon_0\}
\]
has the following properties.  Put
\(\omega_n=v_{n,a}\) and \(d_n=b_{n,a}\) for \(n\in\mathcal N_a\).
Then
\[
 \lvert d_n\rvert
 \leq C\bigl(\tau-\theta_n+\gamma^n\bigr).
\tag{51}\label{eq:d-tail}
\]
For every \(\eta>0\), there are \(C_\eta>0\) and \(N_\eta\geq N\) such
that
\[
 d_n^2\leq\eta e_n+C_\eta\gamma^{2n}
 \qquad(n\in\mathcal N_a,\ n\geq N_\eta).
\tag{52}\label{eq:d-absorbed}
\]
For every \(i\neq a\),
\[
 \sum_{n\in\mathcal N_a}b_{n,i}^2<\infty.
\tag{53}\label{eq:other-errors-summable}
\]
\end{lemma}

\begin{proof}
Equation \Cref{eq:active-error-representation} at \(i=a\), together with
uniform boundedness away from zero of \(B_a\), gives
\[
 \lvert d_n\rvert
 \leq C(\sigma_n+\lvert\delta_n\rvert+\epsilon_n).
\tag{54}\label{eq:d-basic}
\]
By \Cref{lem:outside-tail},
\(\Delta_n+\epsilon_n\leq
C(\tau-\theta_n+\gamma^n)\), and
\(\sigma_n=O(\Delta_n)\).  Equation \Cref{eq:defect-tail} proves
\Cref{eq:d-tail}.

Put \(g_n=\tau-\theta_n\).  By \Cref{eq:sigma-e-relations},
\(e_n\geq g_n/\tau\), while \Cref{eq:d-tail} gives
\[
 d_n^2\leq C(g_n^2+\gamma^{2n}).
\]
Since \(g_n\to0\), for every \(\eta>0\) one has
\(Cg_n^2\leq\eta e_n\) for all sufficiently large \(n\).  This proves
\Cref{eq:d-absorbed}.

For \Cref{eq:other-errors-summable}, let
\[
 \widehat C_{n,i}=\frac{C_n(\xi_i)}{K_n}.
\]
Solving \Cref{eq:active-error-representation} at \(i=a\) and substituting
into the same formula at \(i\neq a\) gives
\[
 b_{n,i}
 =\sigma_n+
 \frac{B_i(\zeta_n)}{B_a(\zeta_n)}
   (d_n-\sigma_n+\widehat C_{n,a})
 -\widehat C_{n,i}.
\tag{55}\label{eq:other-error-exact}
\]
By \Cref{eq:endpoint-distance},
\[
 \left\lvert\frac{B_i(\zeta_n)}{B_a(\zeta_n)}\right\rvert
 \leq C\lvert\zeta_n-\xi_a\rvert
 \leq C(\omega_n+\epsilon_n).
\tag{56}\label{eq:B-ratio}
\]
Since \(\widehat C_{n,i}=O(\epsilon_n)\) and \(d_n\to0\),
\[
 \lvert b_{n,i}\rvert
 \leq C\bigl(\sigma_n+\epsilon_n+\omega_n\lvert d_n\rvert\bigr).
\tag{57}\label{eq:other-error-bound}
\]
Now \(\sigma_n\to0\), and eventually
\[
 \sigma_n\leq2\log(1+\sigma_n)
 =2\log\frac{\theta_{n+1}}{\theta_n}.
\]
Thus \(\sum_n\sigma_n<\infty\) by telescoping.  Also
\(\sum_n\epsilon_n<\infty\), and
\[
 \sum_{n\in\mathcal N_a}\omega_nd_n^2
 \leq\sum_{n\in\mathcal N_a}\Delta_n<\infty.
\]
Squaring \Cref{eq:other-error-bound} and using
\(\omega_n^2d_n^2\leq\omega_nd_n^2\) proves
\Cref{eq:other-errors-summable}.
\end{proof}

\subsection{Exclusion of resonance and four-point collapse}

We now use an exact logarithmic cocycle to exclude every surviving
multiplier-\(-1\) coordinate.  Cubic interpolation cancels all terms that
are linear in the four active multiplier errors.

\begin{lemma}[Exclusion of multiplier-minus-one resonances]
\label{lem:no-surviving-resonance}
If \(\Omega\) contains a four-positive vector, then
\(\mathcal R=\varnothing\).
\end{lemma}

\begin{proof}
Suppose to the contrary that \(\lambda=\lambda_j\) for some
\(j\in\mathcal R\), and put
\[
 \ell_i=\ell_i(\lambda)
 =\frac{L(\lambda)}{D_i(\lambda-\xi_i)}.
\]
Lagrange interpolation gives \(\sum_i\ell_i=1\).  Every active coordinate
is nonzero at every finite time: if one vanished, it would remain zero,
contradicting the existence of a four-positive omega-limit point.  For all
sufficiently large \(n\), \(q_n(\xi_i)>0\) and \(q_n(\lambda)<0\), so
\[
 \Psi_{n,\lambda}
 =\log w_{n,j}+\sum_{i=1}^4\ell_i\log v_{n,i}
\tag{58}\label{eq:Psi}
\]
is well defined.  The exact two-step squared-weight recurrence gives
\[
 \frac{\Psi_{n+2,\lambda}-\Psi_{n,\lambda}}2
 =
 \log\frac{\lvert q_n(\lambda)\rvert}{K_n}
 +\sum_{i=1}^4\ell_i
     \log\frac{q_n(\xi_i)}{K_n}.
\tag{59}\label{eq:Psi-half}
\]

The polynomial \(q_n-q_\star\) has degree at most three.  Interpolating it
at the four active nodes, and using
\(q_\star(\lambda)=-\beta\) and
\(q_\star(\xi_i)=\beta\), yields
\[
 \frac{\lvert q_n(\lambda)\rvert}{K_n}
 =1+2e_n-\sum_{i=1}^4\ell_i b_{n,i}.
\tag{60}\label{eq:resonant-interpolation}
\]
Thus, with
\(\Phi_{n,\lambda}=(\Psi_{n+2,\lambda}-\Psi_{n,\lambda})/2\),
\[
 \Phi_{n,\lambda}
 =
 \log\left(1+2e_n-\sum_i\ell_i b_{n,i}\right)
 +\sum_i\ell_i\log(1+b_{n,i}).
\tag{61}\label{eq:Phi-exact}
\]
Taylor expansion at the origin cancels every term linear in \(b_{n,i}\):
\[
 \Phi_{n,\lambda}
 =2e_n+O\left(e_n^2+\sum_{i=1}^4b_{n,i}^2\right).
\tag{62}\label{eq:Phi-Taylor}
\]
All logarithm arguments tend to one.  Because \(e_n\geq0\) and
\(e_n^2=o(e_n)\), there is a constant \(C_\lambda\) such that
\[
 \Phi_{n,\lambda}
 \geq e_n-C_\lambda\sum_{i=1}^4b_{n,i}^2
\tag{63}\label{eq:Phi-lower}
\]
for all sufficiently large \(n\).

Choose the endpoint neighborhoods from \Cref{lem:endpoint-errors}.  Their
complement in a fixed neighborhood of \([\xi_2,\xi_3]\) is compactly
contained in the middle gap.  During visits to this compact interior, all
four active weights are uniformly positive; otherwise the limiting
corrected barycentric equations would place the cluster value at an
endpoint.  Hence
\[
 \Delta_n\geq\sum_{i=1}^4v_{n,i}b_{n,i}^2
\tag{64}\label{eq:interior-energy}
\]
and \(\sum_n\Delta_n<\infty\) make
\(\sum_i b_{n,i}^2\) summable over the interior visits.

At a visit in \(\mathcal N_a\), the errors \(b_{n,i}^2\) for
\(i\neq a\) are summable by
\Cref{eq:other-errors-summable}.  Apply
\Cref{eq:d-absorbed} with
\(\eta<(2C_\lambda)^{-1}\).  Then
\Cref{eq:Phi-lower} gives
\[
 \Phi_{n,\lambda}
 \geq\frac12e_n
 -C_\lambda\sum_{i\neq a}b_{n,i}^2-C\gamma^{2n}
 \qquad(n\in\mathcal N_a)
\tag{65}\label{eq:Phi-endpoint}
\]
for all sufficiently large \(n\).  Every large index is either an interior
visit or lies in one of the two endpoint neighborhoods, because every
cluster value of \(\zeta_n\) lies in \([\xi_2,\xi_3]\).  It follows that
the negative parts of \(\Phi_{n,\lambda}\) are summable on each parity.
Since \(\Psi_{n+2,\lambda}-\Psi_{n,\lambda}=2\Phi_{n,\lambda}\), this
summability bounds \(\Psi_{n,\lambda}\) below on each parity.

Choose an even subsequence converging to the four-positive omega-limit
point.  Along it, all four active logarithms in \Cref{eq:Psi} remain
bounded, whereas \(w_{n,j}\to0\) by
\Cref{lem:off-support-summability}.  Hence
\(\Psi_{n,\lambda}\to-\infty\) along that subsequence, contradicting the
parity-wise lower bound.  Therefore \(\mathcal R=\varnothing\).
\end{proof}

With resonances absent, every surviving transverse coordinate has a strict
gap.  The resulting geometric forcing and the parabolic dichotomy make the
corrected barycentric coordinate converge.

\begin{lemma}[Strict-gap four-point collapse]
\label{lem:strict-gap-collapse}
Suppose
\[
 \lvert q_\star(\lambda_j)\rvert<\beta
\]
for every \(j\notin I_\star\) whose coordinate survives.  Then
\(\Omega\) is a singleton and \(y_{2k}\) converges.
\end{lemma}

\begin{proof}
If every vector in \(\Omega\) has grade three, the claim follows from
\Cref{lem:three-point-collapse}.  Otherwise, \(\Omega\) contains a
four-positive vector.  No active coordinate can then have been annihilated
at a finite time.

The strict inequalities, full-sequence convergence \(q_n\to q_\star\),
and \(H_n\to\beta\) give constants \(C_0>0\) and \(0<\gamma<1\) such
that
\[
 \epsilon_n\leq C_0\gamma^n
\tag{66}\label{eq:strict-gap-geometric}
\]
for all sufficiently large \(n\).  This follows by applying
\Cref{eq:outside-two-step} separately on the two parities; coordinates
annihilated at a finite time remain zero.

By \Cref{lem:resonant-geometry-tail}, the limiting active sign pattern is
\((+,-,-,+)\), and every cluster value of \(\zeta_n\) belongs to
\([\xi_2,\xi_3]\).  The corrected barycentric equations also give
\[
 \operatorname{dist}(\zeta_n,[\xi_2,\xi_3])
 \leq C\epsilon_n.
\tag{67}\label{eq:zeta-gap-distance}
\]
Indeed, let \(s_i\) be the eventual sign of \(p_n(\xi_i)\).  From
\Cref{eq:corrected-barycentric} and boundedness of the outside correction,
\[
 s_iH_n\frac{\xi_i-\zeta_n}{D_i}\geq-C\epsilon_n
 \qquad(1\leq i\leq4).
\]
The four exact half-line inequalities intersect in
\([\xi_2,\xi_3]\), and their \(O(\epsilon_n)\) relaxations imply
\Cref{eq:zeta-gap-distance}.  By \Cref{eq:X-limit},
\(X(p_{n+1},z)\) is bounded away from zero for large \(n\) and
\(z\) near this interval.

Equation \Cref{eq:full-zeta-increment} and
\Cref{eq:forcing-coefficient-bound} give
\[
 \lvert\zeta_{n+1}-\zeta_n\rvert
 \leq C(\lvert\delta_n\rvert+\epsilon_n).
\]
Use this in \Cref{eq:full-defect-transport}, subtract the two \(X\)
terms, and use their uniform Lipschitz dependence on the second argument.
Since \(\delta_n\to0\),
\[
 \lvert\delta_{n+1}-\delta_n\rvert
 \leq C\delta_n^2+C\gamma^n.
\tag{68}\label{eq:strict-defect-parabolic}
\]
\Cref{lem:parabolic-dichotomy} applies.  If
\(\sum_n\lvert\delta_n\rvert<\infty\), then
\Cref{eq:full-zeta-increment} and \Cref{eq:strict-gap-geometric} give
finite total variation of \(\zeta_n\).

Otherwise \(\delta_n\) is eventually one-signed.  The polynomial \(L\) is
nonnegative on \([\xi_2,\xi_3]\), while
\Cref{eq:zeta-gap-distance} shows that the negative part of
\(L(\zeta_n)\) is \(O(\epsilon_n)\).  After orienting the real line
according to the fixed sign of \(\delta_n\),
\Cref{eq:full-zeta-increment} shows that all increments in the adverse
direction have a summable total, because \(F_n=O(\epsilon_n)\) and
\(\sum_n\epsilon_n<\infty\).  A bounded real sequence with summable
increments in one direction converges: adding the tail of those adverse
increments makes it monotone up to a vanishing correction.  Thus
\(\zeta_n\) converges in both alternatives.

For even \(n=2k\), \Cref{eq:corrected-barycentric} gives
\[
 v_{2k,i}p_{2k}(\xi_i)
 =H_{2k}\frac{\xi_i-\zeta_{2k}}{D_i}+O(\epsilon_{2k}).
\]
Because
\(p_{2k}(\xi_i)\to p_{\mathrm e}(\xi_i)\neq0\), all active squared
coordinates converge, while the outside squared coordinates tend to zero.
Finally,
\[
 \frac{y_{2k+2,\iota_i}}{y_{2k,\iota_i}}
 =\frac{q_{2k}(\xi_i)}{K_{2k}}>0
\]
for every sufficiently large \(k\).  Hence the sign of each nonvanishing
active coordinate eventually stabilizes, and a coordinate whose squared
weight tends to zero itself tends to zero.  Therefore \(y_{2k}\) converges,
and \(\Omega\) is a singleton.
\end{proof}

The preceding two lemmas classify every possible outside coordinate and
complete the four-node case without a genericity or strict-gap assumption.

\begin{lemma}[Collapse of the four-point omega-limit set]
\label{lem:four-point-collapse}
Assume \(\lvert I_\star\rvert=4\).  Then \(\Omega\) is a singleton.
Consequently, \(y_{2k}\) converges and, by continuity of \(T\), so does
\(y_{2k+1}\).
\end{lemma}

\begin{proof}
If every vector in \(\Omega\) has grade three,
\Cref{lem:three-point-collapse} gives the conclusion.  Otherwise,
\(\Omega\) contains a grade-four vector, necessarily four-positive on
\(I_\star\).

Let \(j\notin I_\star\) be a surviving coordinate.  By
\Cref{eq:outside-two-step} and
\Cref{lem:off-support-summability}, the inequality
\(\lvert q_\star(\lambda_j)\rvert>\beta\) is impossible: its squared
two-step multiplier would eventually exceed one, whereas the coordinate
mass tends to zero.  The equality
\(q_\star(\lambda_j)=\beta\) is impossible by the definition of
\(I_\star\), and
\Cref{lem:no-surviving-resonance} excludes
\(q_\star(\lambda_j)=-\beta\).  Thus every surviving outside coordinate
satisfies
\(\lvert q_\star(\lambda_j)\rvert<\beta\); every nonsurviving coordinate
is eventually zero.  \Cref{lem:strict-gap-collapse} applies, so
\(\Omega\) is a singleton.  Its point has grade three or four by
\Cref{lem:omega-structure}, and continuity of \(T\) there gives convergence
of the odd subsequence.
\end{proof}

\subsection{Conclusion}

\begin{proof}[Proof of \Cref{thm:forsythe-s2}]
\Cref{lem:energy-ratio} proves that every residual is nonzero and retains
grade at least three.  By
\Cref{lem:iteration-equivalence,lem:spectral-reduction}, the normalized
residuals are isometric to the diagonal monic-quadratic orbit
\(y_{k+1}=T(y_k)\).

\Cref{lem:monic-monotonicity,lem:omega-structure,%
lem:quartic-convergence} produce a nonempty connected even omega-limit set
\(\Omega\) and an equality set \(I_\star\) with
\(3\leq\lvert I_\star\rvert\leq4\).  If every omega-limit vector has grade
three, \Cref{lem:three-point-collapse} makes \(\Omega\) a singleton.
Otherwise, some omega-limit vector has grade four, so
\(\lvert I_\star\rvert=4\), and
\Cref{lem:four-point-collapse} again makes \(\Omega\) a singleton.
Therefore \(y_{2k}\to y_{\mathrm e}\).  By
\Cref{lem:omega-structure}, \(y_{\mathrm e}\) has grade three or four,
\(T\) is continuous there, and
\[
 y_{2k+1}=T(y_{2k})\longrightarrow
 y_{\mathrm o}:=T(y_{\mathrm e}),\qquad
 T(y_{\mathrm o})=y_{\mathrm e}.
\]
The proof of \Cref{lem:omega-structure} gives grade at least three for
\(y_{\mathrm o}\).  Moreover, with
\(q(t)=\pi_{y_{\mathrm o}}(t)\pi_{y_{\mathrm e}}(t)\), polynomial
commutativity and \(q(A)y_{\mathrm e}=\beta y_{\mathrm e}\) imply
\(q(A)y_{\mathrm o}=\beta y_{\mathrm o}\); hence its grade is at most
four.  The fixed spectral isometry in \Cref{lem:spectral-reduction}
transfers both limits and their grades to the original residual
directions.  The energy-ratio estimate is
\Cref{eq:energy-ratio-bound}, and the per-restart resource bound follows
from \Cref{lem:iteration-equivalence}.  This proves the theorem.
\end{proof}

\bibliographystyle{alpha}
\bibliography{paper}

\end{document}
